\documentclass[a4paper]{lipics-v2021}


\title{Learning Visibly Pushdown Languages under Rewriting Rules}

%\titlerunning{Approximate Learning of Limit-Average Automata}%optional, please use if title is longer than one line

\author{Jan Otop}{University of Wroc\l{}aw}{jan.otop@uwr.edu.pl}{https://orcid.org/0000-0002-8804-8011}{}
\author{Piotr Traczyński}{University of Wroc\l{}aw}{}{}{} %mandatory, please use full name; only 1 author per \author macro; first two parameters are mandatory, other parameters can be empty.

\authorrunning{Jan Otop and Piotr Traczyński}%mandatory. First: Use abbreviated first/middle names. Second (only in severe cases): Use first author plus 'et. al.'

\Copyright{Jan Otop and Piotr Traczyński}%mandatory, please use full first names. LIPIcs license is "CC-BY";  http://creativecommons.org/licenses/by/3.0/


\ccsdesc[500]{Theory of computation~Grammars and context-free languages}
\ccsdesc[500]{Theory of computation~Rewrite systems}
%\ccsdesc[100]{Mathematics of computing~Markov networks}
%}% mandatory: Please choose ACM 2012 classifications from https://www.acm.org/publications/class-2012 or https://dl.acm.org/ccs/ccs_flat.cfm . E.g., cite as "General and reference $\rightarrow$ General literature" or \ccsdesc[100]{General and reference~General literature}. 

\keywords{visibly pushdown automata, string rewriting, active learning} 

\usepackage[dvipsnames]{xcolor}

\usepackage{amsmath, amsthm}
\usepackage{todonotes}
\presetkeys{todonotes}{inline,backgroundcolor=green}{}
\usepackage{amssymb}
\usepackage{xspace}
\usepackage{stmaryrd} % Where \llbracket and \rrbracket are defined


\newcommand{\Paragraph}[1]{\noindent\textbf{#1.}}


%\newtheorem{theorem}{Theorem}
%\newtheorem{lemma}[theorem]{Lemma}
%\newtheorem{definition}[theorem]{Definition}
\newtheorem{fact}[theorem]{Fact}
%\newtheorem{example}[theorem]{Example}
%\newtheorem{proposition}[theorem]{Proposition}
%\newtheorem{remark}[theorem]{Remark}
%\newtheorem{corollary}[theorem]{Corollary}


\newcommand{\N}{\mathbb{N}}
\newcommand{\set}[1]{\{#1\}}

\newcommand{\arity}{\textsf{Ar}}
\newcommand{\dom}{\textrm{dom}}
\newcommand{\countAut}[2]{\#_{#1}[#2]}
\newcommand{\countAutCumm}[2]{\Sigma\#_{#1}[#2]}


% 
\newcommand{\coNP}{\textsf{coNP}}
\newcommand{\NP}{\textsf{NP}}
\newcommand{\PTime}{\textsf{P}}
\newcommand{\PSPACE}{\textsf{PSPACE}}
\newcommand{\PTIME}{\textsf{P}}



% Rewriting
\newcommand{\rewritesOneStep}{\rightarrow}
\newcommand{\rewrites}{\rewritesOneStep^*}
%\newcommand{\normalForm}[1]{\textbf{NF}(#1)}
\newcommand{\normalFormTRS}[2]{\textbf{NF}_{#1}(#2)}
\newcommand{\rew}{\mathcal{R}}

\newcommand{\subst}{\sigma}
\newcommand{\substExt}{\widehat{\subst}}


% Automata

%\newcommand{\Sig}{\mathcal{F}}


\newcommand{\aut}{\mathcal{A}}
\newcommand{\autB}{\mathcal{B}}
\newcommand{\lang}{\mathcal{L}}

\newcommand{\Qfin}{Q_{\textrm{acc}}}

\newcommand{\deltaFinal}{\widehat{\delta}}
%\newcommand{\deltaFinal}{\widehat{\delta}}
\newcommand{\deltaExt}[1]{\vec{\delta}_{#1}}

\newcommand{\congL}{\sim_{\lang}}

% Other

\newcommand{\lStar}{L\ensuremath{^*}}
\newcommand{\Vars}{\mathcal{X}}


\newcommand{\targetLang}{\mathcal{U}}


\newcommand{\subsumedBy}{\leq_{\aut}}
\newcommand{\subsumedByStrict}{<_{\aut}}
\newcommand{\subsumedByFunc}{\preccurlyeq_{\aut}}
\newcommand{\succBy}{\textsf{succ}_{\aut}}


\newcommand{\Iff}{\leftrightarrow}


%% DL
\newcommand{\DLand}{\sqcap}
\newcommand{\DLor}{\sqcup}

\newcommand{\TBD}{{\color{green}{XX}}\xspace}

\newcommand{\Lwords}{\mathcal{K}}
\newcommand{\SigUnary}{\Sig_{\textrm{unary}}}



\newcommand{\lhs}{\boldsymbol{\alpha}}
\newcommand{\rhs}{\boldsymbol{\beta}}
\newcommand{\alphabet}{\Sigma}
\newcommand{\vAlph}{\widehat{\Sigma}}


\newcommand{\DFA}{\textbf{DFA}\xspace}
\newcommand{\DVPA}{\textbf{DVPA}\xspace}
\newcommand{\VPA}{\textbf{VPA}\xspace}
\newcommand{\VPL}{\textbf{VPL}\xspace}
\newcommand{\VPLs}{\textbf{VPLs}\xspace}

\newcommand{\stackAlphabet}{\Gamma}
\newcommand{\stackAlphabetBot}{\Gamma \cup \set{\bot}}
\newcommand{\trans}[2]{\rightarrow_{#1}^{#2}}
\newcommand{\oracle}{\mathcal{O}}

\newcommand{\SRS}{\textbf{SRS}\xspace}
%\newcommand{\SRSV}{\textbf{SRS-V}\xspace}
\newcommand{\SRSs}{\textbf{SRSs}\xspace}
\newcommand{\rewV}{\mathcal{R[\Vars]}}
\newcommand{\sem}[1]{\llbracket #1 \rrbracket}


\newcommand{\bigO}{\mathcal{O}}


\usepackage{thmtools} 
\usepackage{thm-restate}
%\usepackage{minted}
\usetikzlibrary{automata, positioning, arrows.meta,bending}

%\usepackage{dutchcal} % \mathcal: also small letters. \mathbcal = bold ones. 

%\hideLIPIcs
\begin{document}

\maketitle

\begin{abstract}
We present a method for approximating Visibly Pushdown Languages (VPLs) using string rewriting rules. 
We apply this framework to extend Angluin-style active learning for VPLs by incorporating deductive inference from rewrite rules. 
String rewriting rules that encode prior knowledge about the target language and  
the learning algorithm can deductively infer answers to some queries without asking the oracle, 
thereby reducing its query complexity.
We demonstrate the efficacy of this approach through examples of rewrite systems that capture natural properties of structured data
and significantly reduce the number of equivalence queries directed to the oracle. 
Finally, we address the rule synthesis problem,  which is a first step for making active automata learning algorithms explainable.
\end{abstract}


\section{Introduction}
\Paragraph{Active automata learning}
The aim of \emph{active automata learning} is to generate a \emph{deterministic finite automaton (\DFA)} recognizing an unknown regular language 
$\lang$ (target language)
given access to an oracle that answers queries about $\lang$.
%Unlike \emph{passive} learning, where the input dataset is fixed in advance, the algorithm may \emph{actively} query the language.
There are two types of queries: 
\emph{membership} of a word in $\lang$, and 
\emph{equivalence} of the language of a given \DFA $\aut$ with the target language $\lang$,
where the oracle either confirms equivalence or returns a \emph{counterexample} word witnessing the difference between $\lang$ and $\lang(\aut)$.
This framework was proposed in the seminal paper~\cite{angluin1987learning} together with the \lStar-algorithm, 
which enables learning in polynomial time. 
Active automata learning has been intensively researched 
for both efficiency~\cite{TTTalgorithm,Lsharp} as well as extensions to other automata-based models such as
nondeterministic finite automata~\cite{BolligHKL09}, tree automata~\cite{DBLP:conf/dlt/DrewesH03}, nominal automata over infinite alphabets~\cite{MoermanS0KS17} 
and weighted automata~\cite{activeWeightedAutomata}.

\Paragraph{Learning context-free languages}
The active learning framework has been extended to pushdown languages~\cite{Knobe76,Angluin95,GrammaticalInference,DBLP:conf/pldi/Bastani0AL17,DBLP:conf/kbse/KulkarniLS21}, 
where the representations typically considered have been context-free grammars.
However, 
generating context-free grammars from queries is substantially more difficult than generating \DFA, as there is no counterpart to the Myhill-Nerode theorem for pushdown languages. 
Furthermore, the equivalence problem for context-free grammars is undecidable, and 
for deterministic context-free grammars, which are equivalent to deterministic pushdown automata,
the best currently known upper bound on the complexity of deciding equivalence is primitive recursive~\cite{DBLP:conf/icalp/Senizergues97,DBLP:journals/jcss/Jancar21}. 
However, a subclass of context-free languages called \emph{visibly pushdown languages (\VPLs)}~\cite{InputDrivenPDA,DBLP:conf/stoc/AlurM04}, where the stack operations are determined by the input letters, 
admits a Myhill-Nerode-style characterization~\cite{alur2005congruences} and 
language equivalence of deterministic visibly pushdown automata (\DVPA) is decidable in polynomial time~\cite{DBLP:conf/stoc/AlurM04}.

\Paragraph{Learning \VPLs}
While \DVPA share multiple good properties with \DFA such as polynomial-time closure over Boolean operations and subset-based determinization, 
the minimization problem for \DVPA is \NP-complete~\cite{DBLP:journals/lmcs/GauwinMR19}, which
rules out polynomial-time active learning algorithms (unless \PTIME=\NP).
There are different approaches to address this problem including 
restricting the learning problem to a subclasses of \DVPA~\cite{DBLP:conf/concur/KumarMV06},
working only with \DVPA in a special form such as 
SEVPA, MeVPA or CDA~\cite{DBLP:phd/dnb/Isberner15, ValidatingXMLwithVPA,DBLP:journals/pacmpl/JiaT24}, 
which may involve exponential blow-up upon transformation to \DVPA~\cite{alur2005congruences}, or
augmenting the oracle with queries that depend on the structure of the target \DVPA~\cite{mfcs22}.

\Paragraph{Applications of learning \VPLs}
Learning \VPLs has been applied to a wide range of domains, 
including 
inference of specifications for document validation from JSON schemas, 
generation of valid program inputs~\cite{DBLP:conf/pldi/Bastani0AL17,DBLP:journals/pacmpl/JiaT24}, and 
verification and explainability of Recurrent Neural Networks (RNNs)~\cite{DBLP:conf/icgi/BarbotBFHKLN0Y21}. 
In these applications, equivalence queries are typically approximated because exact equivalence checking is 
often infeasible or undecidable. 
For instance, verifying the equivalence of a \DVPA against a general program is undecidable,
 while comparison with an RNN inevitably yields a negative result, as standard RNNs are limited to recognizing regular languages. 
 Furthermore, the non-equivalence witness can be prohibitively large~\cite{DBLP:conf/emnlp/HewittHGLM20}, 
 making the negative answer meaningless. 
These factors make answering equivalence queries a significant challenge in such settings.

\Paragraph{Query complexity}
Due to high cost of answering queries, 
there is a large body of work on reducing the query complexity of active automata learning~\cite{TTTalgorithm,ADTalgorithm,Lsharp,KrugerJR24,DierlFHJST24} as well as implementing the teacher
using Large Language Models~\cite{learningEL,bhattamishraautomata,vazquezchanlatte2025ll}.
However, 
LLMs can only answer membership queries, and equivalence query answers are merely approximated by a large number of membership queries, 
which is also the standard approach to equivalence queries in other settings~\cite{ModelLearning}.
Consequently, we adopt the alternative approach proposed in~\cite{ecai25}, 
which restricts the algorithm's search space by supplying an additional information regarding the target language called \emph{advice}.
Specifically, we focus on inferring responses to equivalence queries, 
as this yields the greatest reduction in query complexity and improves the overall reliability of the learning process, particularly in contexts where equivalence queries are typically approximated.
    
\Paragraph{Learning with advice}
    We consider the active learning framework augmented with advice, 
where the learning algorithm is provided with \emph{advice} expressing prior knowledge regarding the target language~\cite{ecai25}. 
By using the advice to deductively infer query answers, this framework bridges two synthesis paradigms:
deductive synthesis from specifications and inductive synthesis via queries. 
This hybrid approach allows a language to be represented partially through queries and partially through advice expressed as a \emph{string rewriting system (SRS)}.
Since SRSs and automata operate on fundamentally different principles, certain language properties admit far more succinct descriptions as rewriting rules than through query-answer interactions.

\Paragraph{String rewriting rules for pushdown languages}
A String Rewriting System (SRS) consists of rewrite rules $l \rightarrow r$, where $l,r$ are words.
\emph{Rewriting} is the process of iteratively transforming words by replacing the infix that matches the left-hand side of a rule $\lhs \to \rhs$ with its corresponding right-hand side $\rhs$.                 
For instance, commutativity of letters $a,b$ can be expressed with $ab \to ba$. 
However, if $a$ is a call letter and $b$ is a return letter, 
the commutativity rule disrupts the matching between call and return letters. 
We discuss this issue in the paper and define a class of rules, which are compatible with the implicit structure of pushdown languages. 
Furthermore, in the pushdown case, there are properties not expressible by simple rewrite rules such as 
all content between comment tags can be discarded. 
We address this issue by extending rewrite rules with variables. 

\Paragraph{Modelling prior knowledge with string rewriting}
In our framework, an advice SRS induces an equivalence relation on words relative to the target language  $\lang$.
Specifically, if a word $w$ belongs to $\lang$, then every word obtained by rewriting $w$ must also be in $\lang$; symmetrically, rewriting must respect the complement of $\lang$. 
We also consider one-sided variants based on implication stating that 
(1)~words from the language may only be rewritten words in the language, while there are no restrictions for words outside the language, or
(2)~words outside the language may only be rewritten words outside the language.
This allows the SRS to capture high-level structural properties, such as commutativity or idempotence, without necessarily characterizing the learned language completely.

\Paragraph{Contributions}
While active automata learning with advice was introduced in~\cite{ecai25}, extending this framework 
from regular languages to \VPLs and Deterministic Visibly Pushdown Automata (DVPAs) is non-trivial. 
This transition requires addressing structural challenges of pushdown systems, which are absent in the regular case.
The main contributions of presented in this paper are:  
\begin{enumerate}
    \item 
    We define \emph{well-matched} and \emph{matching-preserving} rules that are compatible with the structural properties of pushdown languages.
    Furthermore, we extend these rules with (context) variables in order to capture natural properties of pushdown language.
    \item We develop an algorithm that infers answers to equivalence queries for \DVPA. 
    \item We establish the complexity of the rule synthesis problem
    it is a functional problem, so we show that a related decision problem is \NP-complete.
    \item We demonstrate the practical benefit of the advice mechanism through experiments.
\end{enumerate}
%The detailed proofs of presented results are relegated to the appendix.

\Paragraph{Related work}
There is a large body of work on active learning of context-free grammars~\cite{Knobe76, Shapiro82, Angluin87a, Sakakibara92, DBLP:conf/emnlp/SennhauserB18, GrammaticalInference, DBLP:conf/pldi/Bastani0AL17, DBLP:conf/pldi/BendrissouGZ22, DBLP:conf/kbse/KulkarniLS21, DBLP:conf/sigsoft/WuJYBSPX19}, as well as in learning visibly pushdown grammars and visibly pushdown automata~\cite{DBLP:phd/dnb/Isberner15, mfcs22, ValidatingXMLwithVPA, DBLP:conf/icgi/BarbotBFHKLN0Y21, DBLP:journals/pacmpl/JiaT24, DBLP:journals/corr/abs-2603-11648}.
These works primarily aim to improve learning algorithms for arbitrary input languages by exploring alternative representation models, augmenting the learning framework with additional capabilities 
(e.g., oracles that answer additional types of queries, or assuming the set of non-terminals is known), 
or altering the mode of learning. These works are mostly orthogonal to ours. 
In our work, we restrict the search space;
improving the performance of active learning algorithms by specializing them to subclasses of languages is 
an active research direction for Mealy machines~\cite{graybox-learning, DBLP:conf/dlt/BerthonBPR21, DBLP:conf/fossacs/LabbafGHM23, DBLP:journals/corr/abs-2405-08647, ProductAutomata} and tree automata~\cite{DBLP:conf/fossacs/HeerdtKR021, BjorklundBE17}. 
%In contrast to these works, we do not commit to any fixed subclass of \VPLs; instead, the search space of our algorithm is parameterized by the advice SRS provided in the input.


\section{Preliminaries}
\Paragraph{Pushdown alphabet}
A pushdown alphabet $\vAlph$ consists of three disjoint sets of letters $\vAlph = (\alphabet_c,\alphabet_l, \alphabet_r)$ called, respectively, 
\emph{call}, \emph{local} and \emph{return} letters.

\Paragraph{Well-matched words}
\newcommand{\dCall}{\textbf{c}}
\newcommand{\dReturn}{\textbf{r}}
The set of well-matched words over a pushdown alphabet $\vAlph$ is the least set containing $\epsilon$ 
such that if $w_1, w_2$ are well-matched, then
(1)~the concatenation $w_1 w_2$ is well-matched, 
(2)~for every call letter $c$ and return letter $r$, the word $ c w r$ is well-matched, and
(3)~for every local letter $l$, words $lw$ and $wl$ are well-matched.
In a word $w$ (not necessarily well-matched), if $w = x \dCall y \dReturn z$ such that $y$ is a well-matched word, 
we say that letters $\dCall$ and $\dReturn$ are matched, and $\dReturn$ is the return letter corresponding to call $\dCall$ and vice versa. 
Formally, positions $|x|+1$ and $|x|+|y|+2$ in $w$ are in the \emph{matching relation}.

\Paragraph{Visibly pushdown automata} 
A \emph{(non-deterministic) visibly pushdown automaton} (\VPA) $\aut$ is a tuple $\aut = (\vAlph, \stackAlphabet, Q, Q_0, \delta, F)$  consisting of:
(1)~$\vAlph$, a finite pushdown alphabet, 
(2)~$\stackAlphabet$, a finite stack alphabet, 
(3)~$Q$, a finite set of states, 
(4)~$Q_0 \subseteq Q$, a set of initial states, 
(5)~$\delta$, a transition relation defined as below, and
(6)~$F \subseteq Q$, a set of accepting states.
The transition relation $\delta$ is the union of relations $\delta_c \cup \delta_l \cup \delta_r$, where
$\delta_c \subseteq Q \times \alphabet_c \times Q \times \stackAlphabet$, 
$\delta_l \subseteq Q \times \alphabet_l \times Q$, and
$\delta_r \subseteq Q \times \alphabet_r \times \stackAlphabetBot \times Q$.

\Paragraph{Deterministic visibly pushdown automata} 
A deterministic \VPA (called $\DVPA$) is a \VPA $(\vAlph, \stackAlphabet, Q, Q_0, \delta, F)$ such that $|Q_0| = 1$ and 
its transition relation $\delta$ is a function, which is the union of three functions:
$\delta_c \colon Q \times \alphabet_c \to Q \times \stackAlphabet$, 
$\delta_l \colon Q \times \alphabet_l \to Q$,
$\delta_r \colon Q \times \alphabet_r \times \stackAlphabetBot \to Q$.

\Paragraph{Semantics of \VPA} 
Consider a \VPA $\aut = (\vAlph, \stackAlphabet, Q, Q_0, \delta, F)$.
 A~\emph{configuration} of $\aut$ is a pair consisting of the current state $q\in Q$ of $\aut$ and 
its current stack content $u\in \bot \cdot \stackAlphabet^*$. 
A~\VPA defines an infinite-state transition system over its configurations $Q \times (\bot \cdot \stackAlphabet^*)$, 
where configurations $Q_0 \times \set{\bot}$ are initial.
For $a \in \vAlph$, let $\trans{\aut}{a}$ be an auxiliary relation defined as follows:
\begin{itemize}
\item $(q,u) \trans{\aut}{a} (q',uB)$, if $a$ is a call letter and $(q,a,q',B) \in \delta_c$, 
%where $q, q'$ are states, $uB$, $u v$ are stack contents with $B$ being a single (top) symbol.
\item $(q,u) \trans{\aut}{a} (q',u)$, if $a$ is a local letter and $(q,a,q') \in \delta_l$, 
\item $(q,uB) \trans{\aut}{a} (q',u)$, if $a$ is a return letter $(q,a,B,q') \in \delta_r$ and $B \in \stackAlphabet$, and
\item $(q,\bot) \trans{\aut}{a} (q',\bot)$, if $a$ is a return letter $(q,a,\bot,q') \in \delta_r$.
\end{itemize}
For a word $w = a_1 \ldots a_n$, we define 
$\trans{\aut}{w}$ as the composition of relations $\trans{\aut}{a_1}, \dots, \trans{\aut}{a_n}$, i.e.,
$(q,u_0) \trans{\aut}{w} (q',u_n)$ if there exist intermediate configurations
$(q_1, u_1), \ldots, (q_{n-1}, u_{n-1})$ such that for each $1 \leq i \leq n$ we have $(q_{i-1}, u_{i-1}) \trans{\aut}{a_i}  (q_{i}, u_{i})$.
A word $w$ is \emph{accepted} by $\aut$ if there are $q_0 \in Q_0$, $q \in F$ and $u \in \bot \cdot \stackAlphabet^*$ such that $(q_0, \bot) \trans{\aut}{w} (q,u)$.
 The \emph{language recognized by $\aut$}, denoted $\lang(\aut)$, is the set of words accepted  by $\aut$. 
%For a \VPA $\aut$, we define the stack content upon reading $w$, 
%denoted by $\stack{\aut}{w}$, as the unique $u$ such that 
%$(q_0, \bot) \trans{\aut}{w} (q',u)$ for some $q'$. 

\Paragraph{String rewriting} 
A \emph{string rewriting system (\SRS)} $\rew$ over an alphabet $\Sigma$ is 
a (possibly infinite) set of pairs of words $(l,r)$ over $\Sigma$.
A pair of words  $(l,r)$  from $\rew$ is called a \emph{string rewrite rule} and denoted with $l \rightarrow r$.
For an \SRS $\rew$, we define a \emph{single-step rewrite relation} $\rewritesOneStep_{\rew}$ over words from $\Sigma^*$ as follows:
for all $s,t \in \Sigma^*$, we have $s \rewritesOneStep_{\rew} t$ if and only if there are words $x,y \in \Sigma^*$ and a rewrite rule
$(l,r) \in \rew$ such that $s = x l y$ and $t = x r y$. 
The (string) rewriting relation $\rewrites_{\rew}$ is the transitive and reflexive closure of $\rewritesOneStep_{\rew}$.
We will omit the $\rew$ subscript if the \SRS is clear from the context. 
For further details and notions on the topic of string rewriting see~\cite{book1993string}.

\Paragraph{Active learning of languages and automata}
The framework of active automata learning assumes an oracle, called a \emph{Minimally Adequate Teacher}, which answers two types of queries about the target language $\lang$: 
\begin{itemize}
\item \textbf{membership queries}: given a $w \in \Sigma^*$, is $w \in \lang$?, and
\item \textbf{equivalence queries}: given an automaton $\aut$, is $\lang(\aut) = \lang$? 
If not, the teacher returns a counterexample, which is a word from the symmetric difference of $\lang(\aut)$ and $\lang$.
\end{itemize}
Additional types of queries are also possible~\cite{mfcs22}.
The goal of the active learning algorithm is to return a minimal-size automaton recognizing the target language. 
The automaton model is connected to the type of the learned language: 
regular languages are represented by \DFA, while \VPLs are represented in active learning by \DVPA, Single-Entry Visibly Pushdown Automata (SEVPA) or 
other subclasses of \DVPA~\cite{alur2005congruences,DBLP:conf/mfcs/ChervetW07}.
Our method focuses on inferring answers to equivalence queries; hence, it does not depend on the specific active learning algorithm used, as long as it asks such queries.




\section{Rewriting strings in visibly pushdown languages}
\newcommand{\BalReach}{\mathcal{B}_{(4)}}
\newcommand{\autPlus}{\aut_{\sigma}}

We adapt an advice mechanism introduced for \DFA in~\cite{ecai25} to \DVPA. The main aim of advice is to 
provide the learning with prior knowledge about the learned language, which is used to reduce the number of equivalence queries.

While extending the framework from regular to visibly pushdown languages is straightforward, 
both frameworks differ in considered rewriting rules.
We characterize which rewrite rules are compatible with visibly pushdown languages and analyze the computational complexity of resolving queries using this advice.

\subsection{Connecting SRS with VPL}
%We adapt the framework of active learning with advice for \DFA~\cite{ecai25} to \DVPA. 
The advice for the learning  algorithm  is given via string rewriting systems (SRSs), which
are related to visibly pushdown languages through the notions of consistency:

\begin{definition}[\cite{ecai25}]
\label{def:consistent}
Let $\vAlph$ be a pushdown alphabet, $\rew$ be an SRS over $\Sig$ and $\lang$ be a visibly pushdown language over $\vAlph$. 
We say that 
\begin{itemize}
\item $\rew$ is \emph{(fully) consistent} with $\lang$ if and only if
for all words $x,y$, if $x \rewrites_{\rew} y$, then $x \in \lang \iff y \in \lang$.  
\item $\rew$ is \emph{positively consistent} with $\lang$ if and only if for all words $x,y$, if $x \rewrites_{\rew} y$, then $x \in \lang \implies y \in \lang$.  
\item $\rew$ is \emph{negatively consistent} with $\lang$ if and only if for all words $x,y$, if $x \rewrites_{\rew} y$, then $x \notin \lang \implies y \notin \lang$.  
\end{itemize}
\end{definition}

Note that an SRS $\rew$ is fully consistent with $\lang$ if and only if it is both positively and negatively consistent.
%While more granular than standard consistency, positive and negative consistency are harder to infer.

Having the notions of consistency present the framework of active learning with advice.

\Paragraph{Learning with advice}
The active \DVPA learning with advice problem is as follows.
Given an \emph{oracle} $\oracle$ answering membership and equivalence queries about the target visibly pushdown language $\lang$ (and possibly other queries), and three \emph{advice} SRSs 
$\rew_{=}$, which is consistent with $\targetLang$, and 
$\rew_{+}$ (resp. $\rew_{-}$), which is positively (resp., negatively) consistent with $\targetLang$, the task is to construct 
a \DVPA recognizing $\lang$. 
The advice is used to approximate the equivalence in order to reduce the number of equivalence queries. 
Therefore, it can be integrated with any framework for active learning of visibly pushdown languages~\cite{DBLP:conf/icgi/BarbotBFHKLN0Y21,DBLP:journals/pacmpl/JiaT24} or \DVPA~\cite{mfcs22} as long as the oracle answers membership and equivalence queries.
In particular, the oracle can answer other types of queries as in~\cite{mfcs22}.


\Paragraph{Reducing the query complexity regarding equivalence oracles}
We can reduce the number of equivalence queries without affecting the learning algorithm.
When the learning algorithm attempts to asks the oracle an equivalence query, we can capture the query and first check whether an advice SRS $\rew$ is consistent (resp., positively consistent or negatively consistent) with
the language of a given candidate automaton $\aut$. 
If it is inconsistent, we know that $\lang(\aut)$ differs from the target language, with which $\rew$ is consistent (resp., positively consistent or negatively consistent).
It suffices to compute a pair of words violating consistency (resp., positively consistent or negatively consistent), i.e., 
words $x,y$ such that $x \rewrites_{\rew} y$ but $x \in \lang \implies y \in \lang$ does not hold.
Then, one of the words $x, y$ has to be a counterexample to $\lang(\aut)$  being equal to the target language. 
It can be decided which one with a single membership query.


\subsection{Rewrite rules for visibly pushdown languages}

While the notion of consistency with the language is the same for \DVPA as for \DFA, the main difference lies in the rewrite rules. 

In the case of regular languages, any pair of words is a valid string rewrite rule, but words in visibly pushdown languages 
have an additional implicit structure (which has been made explicit in~\cite{DBLP:journals/jacm/AlurM09}).
The intuition behind rewrite rules is that words are rewritten to other words that are equivalent w.r.t. a given language. 
Therefore, rewrite rules should preserve the matching of pushdown letters except for letters in subwords being rewritten.
Intuitively, rewriting one part of the word, should not affect other parts of the word, which is trivially true for
words, but not for pushdown words. For instance, by rewriting $c_1 c_2 c_3 r_3 r_2 r_1$ with the rule $c_3 r_3 \to \epsilon r_3 c_3$
we get a well-matched word with a very different matching $c_1 c_2 r_2 c_3 r_2 r_1$.
%the matching $(c_1, r_1), (c_2, r_2), (c_3, r_3)$ changes to $(c_1, r_1), (c_2, r_3), (c_3 r_2)$

\begin{definition}
A rewrite rule $\lhs \to \rhs$ \emph{preserves matching} over $\vAlph$ if and only if for every word $w \in \vAlph^*$]
such that $w = x \lhs y$ for some $x,y$ we have:
for every call letter $c$ from $x$ in $\lhs$ has the corresponding return in $y$ 
if and only if the same letter $x$ in  $w' = x \rhs y$ has the corresponding return letter in $y$.
\end{definition}

Matching preserving rules does not brake matching between calls and returns that are not in the rewritten fragment of the word.
A prominent example of matching preserving rules are \emph{well-matched rules}, which are natural and enjoy more efficient algorithms 
for checking consistency:

\begin{definition}
A rewrite rule $\lhs \to \rhs$ is well-matched over $\vAlph$ if and only if words $\lhs$ and $\rhs$ are well-matched.
\end{definition}

Clearly, every well-matched rule preserves matching. However, a rewrite rule $c_1 c_2 r_2 \to c_1$ is a matching preserving rule, while it
is not a well-matched rule. Still, most of our examples are well-matched rules.

While we have restricted the set of admissible rules set to affect pushdown words locally, 
we must also extend it to capture the inherent properties of pushdown languages. 
We demonstrate this in the example below.

\begin{example}
\newcommand{\OpenDelOne}{\ensuremath{\boldsymbol{\langle}}\,}
\newcommand{\CloseDelOne}{\ensuremath{\boldsymbol{\rangle}}\,}
\newcommand{\OpenDelTwo}{\ensuremath{\boldsymbol{[}}\,}
\newcommand{\CloseDelTwo}{\ensuremath{\boldsymbol{]}}\,}
Consider the Dyck language $\lang_2$ of well-matched brackets over tow pairs of brackets: \( \set{ \OpenDelOne, \CloseDelOne} \) and \( \set{\OpenDelTwo, \CloseDelTwo} \). 
The language $\lang_2$ can be approximated by the property that the order of \OpenDelOne and \OpenDelTwo is irrelevant as long as each opening bracket matches the closing one. 
For instance, in every word $\OpenDelOne \OpenDelTwo \CloseDelTwo \CloseDelOne \to  \OpenDelTwo  \OpenDelOne \CloseDelOne \CloseDelTwo$. 
Unfortunately, this property cannot be expressed by a finite set of string rewriting rules as it requires  far apart letters to switch; for example 
$\OpenDelOne \OpenDelTwo^n \OpenDelOne \CloseDelOne \CloseDelTwo^n \CloseDelOne$ is equivalent to $\OpenDelOne \OpenDelOne \OpenDelTwo^n \CloseDelTwo^n \CloseDelOne \CloseDelOne$. 
However, we can express this property with a word variable $X$, for which any well-matched word can be substituted. Then, the rule
\[
\OpenDelOne \OpenDelTwo X \CloseDelTwo \CloseDelOne \to  \OpenDelTwo  \OpenDelOne X \CloseDelOne \CloseDelTwo
\]
expresses that the relative order of \OpenDelOne and \OpenDelTwo is irrelevant as long as the corresponding brackets match.
\end{example}

\todo{Is consistency defined for rules with variables?}
\Paragraph{String rewriting systems with variables}
Consider an infinite set of variables $\Vars$. 
A string rewriting system with variables (\SRSV) $\rewV$ is a finite set of pairs of words $(\lhs,\rhs)$ over $\vAlph \cup \Vars$ such that 
for every rule $\lhs \to \rhs \in \rewV$, 
(1)~in each $\lhs$ and $\rhs$, every variable $X \in \Vars$ occurs at most once, and  
(2)~if a variable $X \in \Vars$ occurs in $\rhs$, then it occurs in $\lhs$ as well.
We consider \SRSVs as finite representations of infinite SRSs. 
That is, we say that a rule $\lhs \to \rhs \in \rewV$ generates an infinite 
SRS $\sem{\lhs \to \rhs}$ obtained by all instantiations of variables from $\lhs$ and $\rhs$ with well-matched words, where variable in both $\lhs$ and $\rhs$ are instantiated to the same word.
E.g. a rule $c X r \to X$ generates an infinite set $\{ c r \to \epsilon,c c r r \to \epsilon cr, \ldots \}$.
Finally, an \SRSV $\rewV$ generates an SRS $\sem{\rewV}$ defined as the union of $\sem{\lhs \to \rhs}$  over all rules $\lhs \to \rhs$ from $\rewV$.

\Paragraph{Matching preserving and well-matched rules with variables}
Consider a rewrite rule with variables $\lhs \to \rhs$. 
We say that $\lhs \to \rhs$ preserves matching (resp. is well-matched) if and only if the rule resulting from deletion of all variables from $\lhs \to \rhs$ preserves matching (resp. is well-matched).
Observe that if $\lhs \to \rhs$ preserves matching (resp. is well-matched), then every rewrite rule for the generated SRS $\sem{\lhs \to \rhs}$ preserves matching (resp. is well-matched).

%Since \SRSV is only a finite representation of infinite \SRS, all notions for \SRSs naturally extend to \SRSVs. 
We only assign well-matched words to variables as otherwise the resulting rules would not preserve matching. 
Consider a rule $c X r \to X$, which is well-matched if it is instantiated with a well-matched word for $X$. 
However, instantiating $X$ with $r c$ results in a rule $c r c r \to r c$, which does not even preserve matching as $c r c r$ is well-matched, while
$r c$ is not.
Furthermore, considered only well-matched instantiations facilitates deciding consistency with a given \DVPA.

\subsection{Deciding consistency}

We discuss the complexity of deciding consistency (resp., positive or negative consistency) of an \SRS or \SRSV with the language of a given \DVPA $\aut$.
Observe that an \SRSV $\rew$ is consistent with the language of $\aut$ if and only if for every rule $\lhs \to \rhs$ and all words $x,y$ we have
$x \lhs y \in \lang(\aut) \iff x \rhs y \in \lang(\aut)$. Indeed, if this property holds for every rule, then by transitivity of the equivalence 
if $w \rewrites_{\rew} w'$, then $w \in \lang \implies w' \in \lang$. Similarly, it suffices to check consistency of single-step rewriting for positive and negative consistency.

\Paragraph{Ground well-matched rules.} Consider a well-matched rule $\lhs \to \rhs$ without variables, i.e., a standard string rewriting rule. 
We extend the alphabet $\vAlph$ with a fresh local letter $\#$ and construct two \DVPA over the extended alphabet.
Both \DVPA extend $\aut$ and differ only in behavior on the letter $\#$; in $\aut_1$ for any state $q$, the transition over $\#$ leads to the state $q'$ such that
$\aut$ moves from $q$ to $q'$ over $\lhs$. Since $\lhs$ is well-matched, the run over $\lhs$ does not depend on the stack content and it does not modify the stack content. 
Therefore, the \DVPA $\aut$ acts of $\lhs$ as on a local letter. 
Next, the \DVPA $\aut_2$ is obtained in a similar way, but $\#$ acts as the word $\rhs$, which is also well-matched. 
Therefore, to check consistency of  $\lhs \to \rhs$  with $\lang(\aut)$ it suffices to check whether there is a word $w$ such that exactly one of \DVPA $\aut_1, \aut_2$ accepts it, 
which can be done by a single reachability check on the product \DVPA $\aut_1 \times \aut_2$, which simulates two copies of $\aut$ with two stacks.
However, the visibility condition ensures that two stacks can be simulated with one over the product of stack alphabets. 
Checking positive (resp., negative) consistency amounts to checking the existence of a word accepted by $\aut_1$ but not $\aut_2$ (resp., $\aut_2$ but not $\aut_1$ for negative consistency), 
which also can be done with a single reachability check. 

\begin{proposition}
Consistency (resp.,  positive, negative consistency) of a ground well-matched rule with the language of a given \DVPA $\aut$ can be decided in polynomial time via a single 
reachability check on a \DVPA obtained from the product of $\aut$ with itself. 
\end{proposition}


Reachability in pushdown automata can be checked in $\bigO(|\aut|^3)$ with a slightly better bound $\bigO(|\aut|^3/\log(|\aut|))$ for \emph{recursive state machines}~\cite{DBLP:conf/popl/Chaudhuri08}, 
which is still $\widetilde{\bigO}(|\aut|^3)$.
Due to the lack of better upper bounds~\cite{DBLP:conf/popl/Chaudhuri08,DBLP:conf/lics/BalasubramanianCM25}, the complexity of reachability in pushdown systems is often referred to as the "cubic bottleneck." 
Consequently, reachability in the product $\aut \times \aut$ amounts to $\widetilde{\bigO}(|\aut|^6)$. 
In the case of \DFA, however, checking consistency of an SRS can be performed locally on the structure of the \DFA without checking reachability. 
The following example illustrates that for pushdown systems, checking consistency must involve reachability.

\begin{example}
Consider a \DVPA $\autB$ depicted in Figure~\ref{fig:nonsyntactic}, which is over the pushdown alphabet $(\set{c},\set{l_1, l_2}, \set{r})$ such that local letters $l_1, l_2$ induce self-loops over all states except for $q_1$. Its stack alphabet is a singleton $\set{a}$.
In the state $q_1$ the \DVPA moves over $l_1$ to $q_2$ and over $l_2$ to $q_2$. The states $q_2, q_3$ have different behavior over the empty stack only, while they are indistinguishable 
over a non-empty stack. Finally, $q_1$ is reachable only with non-empty stack, i.e., $q_1$ is reachable, but the configuration $(q_1, \bot)$ is not.
Then, the rewrite rule $l_1 \to l_2$ is consistent with $\lang(\autB)$, but it requires reachability analysis for that. 
\begin{figure}[h!]
\centering
\begin{tikzpicture}[
    shorten >=1pt, 
    node distance=3.0cm, 
    on grid, 
    auto,
    >={Stealth[bend]}
]

  % State Definitions
  \node[state] (q0) {$q_0$};
  \node[state] (q1) [right=of q0] {$q_1$};
  \node[state] (q2U) [above right=of q1] {$q_2$};
  \node[state] (q2D) [below right=of q1] {$q_3$};
  \node[state] (sink) [below right=of q2U] {$q_R$};
  \node[state,accepting] (acc) [right=of q2U] {$q_F$};

  % Transition Edges
  \path[->]
    (q0) edge node [above] {$c, push(a)$} (q1)
    (q1) edge node [above] {$l_1$} (q2U)
    (q1)  edge node [above] {$l_2$} (q2D)
    (q2U) edge node [right] {$r, pop(a)$} (sink)
    (q2U) edge node [above] {$r, pop(\bot)$} (acc)
    (q2D)  edge [bend right=20] node [above left=0.6cm] {$r, pop(a)$} (sink) 
    (q2D)  edge [bend left=20] node [below right=0.6cm] {$r, pop(\bot)$} (sink);
    

    %[bend right=45]
 %   (q1)  edge [bend right=45] node [above] {$\Sigma\setminus (Y \cup \set{x})$} (qn)
    % The outgoing edge xi from q1
%    (q1)  edge [nodes={at={(1, -0.5)}}] node {$x$} ++(1.5, 0);
%    (q1)  edge node [above] {$x$} (qF)
%    (q1) edge node [left] {$Y$} (sink)
%    (q2) edge node [left] {$Y$} (sink)
%    (qn) edge node [above] {$Y$} (sink)
%    (qF) edge node [left] {$Y$} (sink)
    ;

 \draw[loop right,->] (acc)  to node[right] {$\vAlph$} (acc);
 \draw[loop right,->] (sink) to node[right] {$\vAlph$} (sink);

\end{tikzpicture}
\caption{The automaton $\autB$. All transitions other than depicted are self-loops}
\label{fig:nonsyntactic}
\end{figure}
\end{example} 


\Paragraph{Non-ground well-matched rules}
\newcommand{\autCon}{\aut_{\textrm{con}}}
\newcommand{\DeltaRule}{\Delta[\lhs \to \rhs]}
Consider a well-matched rule $\lhs \to \rhs$ with variables. 
The general approach to decide consistency is, as in the ground case, to reduce it to reachability in a \VPA $\autCon$.
This automaton simulates two copies of $\aut$ over the alphabet extended with a fresh local letter $\#$.
The \VPA $\autCon$ is based on the product $\aut \times \aut$, where the stack simulates the stacks of both copies of the \DVPA $\aut$.
The fresh local letter $\#$ simulates the word $\lhs$ in the first copy and the word $\rhs$ in the second copy.
However, there are two main differences from the ground case.
First, as $\lhs$ and $\rhs$ can be instantiated to (infinitely many) different words, transitions over $\#$ are non-deterministic.
Second, the transitions in each copy of $\aut$ cannot be defined independently; they need to be define on the product to ensure that $\lhs$ and $\rhs$ are instantiated to the same word. 
E.g., for a rule $c X r \to X$, there may be a transition in $\aut$ from $q$ to $q_1$ over $c w_1 r$ and from $q$ to $q_2$ over $w_2$, but we must ensure that $w_1 = w_2$.
Therefore, we compute a set $\DeltaRule$ triples $(q, s_1, s_2)$ such that there are $\alpha, \beta$ satisfying:
(1)~$(q,u) \trans{\aut}{\alpha} (s_1,u)$ and $(q,u) \trans{\aut}{\beta} (s_2,u)$ for every stack content $u$, and
(2)~$\alpha \to \beta$ is an instance  of $\lhs \to \rhs$ obtained by substituting well-matched words for variables.
Using $\DeltaRule$, we define a \VPA $\autCon$ over the pushdown alphabet $\vAlph$ extended with a local letter $\#$ as an extension of $\aut \times \aut$ such that for a state $(q_1, q_2)$ of $\autCon$:
(a)~if $q_1 \neq q_2$, then $\autCon$ has no transition  over $\#$, and
(b)~if $q_1 = q_2$, then for every $(q_1, s_1, s_2) \in \DeltaRule$, 
the tuple $\langle (q_1, q_2), \#, (s_1, s_2) \rangle$ is a transition in $\autCon$. 
Finally, the automaton accepts if in the reached state $(q_1, q_2)$, exactly one of $q_1, q_2$ is accepting in $\aut$. 
Observe that the \VPA $\autCon$ accepts exactly words $w$ that are witnesses of inconsistency of $\lhs \to \rhs$ with $\lang(\aut)$.
These are words $w$ with exactly one occurrence of $\#$ such that there is an instance $\alpha \to \beta$ of $\lhs \to \rhs$ obtained by substituting well-matched word for variables, and
by substituting $\#$ with $\alpha$ and $\beta$ we obtain words $w_1, w_2$ respectively satisfying 
(i)~$w_1 \to w_2$, and 
(ii)~exactly one of $w_1, w_2$ belongs to $\lang(\aut)$.
By modifying the acceptance condition, we can also decide positive and negative consistency. 

In the following, we focus on computing $\DeltaRule$ for rules with variables. 
It suffices to compute how well-matched words transform pairs of states of the \DVPA.

\Paragraph{Computing $\DeltaRule$}
We compute a quaternary relation $\BalReach \subseteq Q^4$, where $(q_1, q_2, s_1, s_2) \in \BalReach$ if and only if there is a well-matched word $w$ such that
$q_1 \trans{\aut}{w} s_1$ and $q_2 \trans{\aut}{w} s_2$. 
It is computed iteratively in a similar way to reachability in pushdown automata.  
\begin{lemma}
Given a \DVPA $\aut$, the relation $\BalReach$ can be computed in polynomial time in the size of $\aut$.
\end{lemma}
Now, we extend the alphabet $\vAlph$ to $\vAlph_+$ with a pair of fresh local letters $x^L$ and $x^R$ for every variable $X$ from $\lhs \to \rhs$.
Next, we substitute each variable $X$ with its corresponding local letters: $x^L$ in $\lhs$ and $x^R$ in $\rhs$ and obtaining words $\alpha^+$ and $\beta^+$ over the extended alphabet. 

We compute $\BalReach$ and then we iterate over all possible assignments from variables of $\lhs \to \rhs$ to $\BalReach$. 
For a given assignment $\sigma$, we extend the \DVPA $\aut$ to a \VPA $\autPlus$ working over the extended alphabet.
If $\sigma(X) = (q_1, q_2, q_1', q_2')$, we define the transitions for the new letters as follows: 
$\autPlus$ has a transition over $x^L$ from $q_1$ to $q_1'$ and over $x^R$ from $q_2$ to $q_2'$; there are no other transitions over $x^L$ or $x^R$.
Finally, we evaluate $\lhs$ and $\rhs$ by evaluating $\alpha^+$ and $\beta^+$ in $\autPlus$.
If $(q,\bot) \trans{\autPlus}{\alpha^+} (s_1,\bot)$ and $(q,\bot) \trans{\autPlus}{\beta^+} (s_2,\bot)$, we add the triple $(q,s_1, s_2)$  to $\DeltaRule$.
Not that the stack content $\bot$ does not influence reachability since $\alpha^+, \beta^+$ are well-matched.


\begin{lemma}
Deciding membership to $\DeltaRule$ over well-matched rules is in \NP.
It can be done in polynomial time assuming a bound on the number of variables in rules.
\end{lemma}

In consequence, checking if there is an inconsistency an be done in \NP and hence:

\begin{theorem}
Consistency (resp., positive, negative consistency) of a well-matched rule with the language of a given \DVPA $\aut$ can be decided in $\coNP$.
Assuming a bound on the number of variables in rules, consistency (resp., positive, negative consistency) of a well-matched rule
can be check in polynomial time.
\end{theorem}


\Paragraph{Complexity of checking (positive) consistency}
The proposed algorithm checks all assignments of transitions for variables, which is exponential in the number of variables, but it works in polynomial time if the number of variables is bounded. 
Improving the proposed exponential-time algorithm is unlikely as the membership problem for $\DeltaRule$ is \NP-complete in general. 
The hardness result hols already for \DFA, which can be considered as \DVPA over only local letters.
Furthermore, the problem of deciding positive (resp. negative) consistency is \coNP-complete as well.

\begin{restatable}{theorem}{NPHardness}
(1)~The problem, given a rewrite rule $\lhs \to \rhs$, a DFA $\aut$ and its states $q, s_1, s_2$, does the tuple $(q, s_1, s_2)$ belongs to $\DeltaRule$, is \NP-complete. 
(2)~The positive (resp. negative) consistency problem over well-matched rules against a regular language given by a \DFA is 
\coNP-complete.
\end{restatable}

The above \coNP-hardness result does not extend to consistency and hence 
the complexity lower bound for checking consistency  of well-matched rules remains an open problem.

Now, we discuss matching-preserving rules, which are not well-matched. 
The overall approach is similar to the well-matched; we extend the alphabet with auxiliary letters.
In contrast to well-matched rules, the new letters are call and return letters. 

\Paragraph{Ground matching-preserving rules}
\newcommand{\autExt}{\aut_E}
\newcommand{\return}{\bar{r}}
\newcommand{\call}{\bar{c}}
Consider a matching preserving rule $\lhs \to \rhs$ without variables, which is not a well-matched rule.
Observe that $\lhs$ and $\rhs$ have the same number of unmatched calls $m$ and unmatched returns $k$, and since it 
is not well-matched, there are unmatched letters (at least one of $k, m$ is non-zero). 
Furthermore,  $\lhs$ and $\rhs$ can be partitioned into well-matched words and unmatched calls and returns, where
unmatched returns precede unmatched calls.
It follows that  
\begin{align*}
\lhs &= v_0 r_1 v_1 \ldots r_k v_k c_1 v_{k+1} c_2 \ldots c_m v_{k+m} \\
\rhs &= v_0' r_1' v_1' \ldots r_k' v_k' c_1' v_{k+1}' c_2' \ldots c_m' v_{k+m}'
\end{align*}
where $v_i, v_i'$ are well-matched word, $r_i, r_i'$ are return letters and $c_i,c_i'$ are call letters. 
Now, as in the well-matched case,  consider a \DVPA $\autExt$ being the product of $\aut$ with itself, i.e., 
the set of states is $Q \times Q$ and the stack alphabet is $\stackAlphabet \times \stackAlphabet$.
The automaton $\autExt$ works over an extended alphabet, where 
the additional letters are $k$ fresh return letters
$\return_1, \ldots, \return_k$
and $m$ fresh call letters $\call_1, \ldots, \call_m$.
The transitions over the letter $\return_1$ simulates the word $v_0 r_1 v_1$  on the first component of $\autExt$
and it simulates the word $v_0' r_1' v_1'$ on the second component.
While doing so, it pops a single letter from the stack, which is a pair of stack symbols of $\aut$. 
Formally, $\delta_E((q_1, q_2),\return_i, (b_1, b_2)) = (s_1, s_2)$ if $(q_1,u b_1) \trans{\aut}{v_0 r_1 v_1} (s_1, u)$ and 
$(q_2,u b_2) \trans{\aut}{v_0' r_1' v_1'} (s_2, u)$ for every stack content $u$. 
For the remaining return letters $\return_i$, where $i \in \set{2, \ldots,k}$
the transitions on $\return_i$ are equivalent to processing  $r_i v_i$  on the first component and $r_i' v_i'$ on the second component.
For new call letter $\call_i$ the transition over $\call_i$ is equivalent to 
the transition over the word $c_i v_{k+1}$ on the first component and $c_i' v_{i+k}'$ on the second component.
Finally, a state $(q_1, q_2)$ of $\autExt$ are accepting if exactly one of $q_1, q_2$ is accepting in $\aut$.

Observe that a word  $w_1 \return_1 \ldots \return_k \call_1 \cdots \call_m w_2$, where $w_1, w_2$ are over the original alphabet $\vAlph$, 
is accepted by $\autExt$  
if and only if $\aut$ accepts exactly one of the words $w_1 \lhs w_2, w_1 \rhs w_2$. 
The latter is equivalent to $\lhs \to \rhs$ being not consistent with $\lang(\aut)$.
Therefore, the rule $\lhs \to \rhs$ is consistent with $\lang(\aut)$ if and only if
the language $\lang(\autExt) \cap \lang(\vAlph^* \return_1 \ldots \return_k \call_1 \cdots \call_m \vAlph^*)$ is empty.

By changing the accepting states of $\autExt$ as discussed in the well-matched case, we can decide positive and negative consistency 
using emptiness of $\lang(\autExt) \cap \lang(\vAlph^* \return_1 \ldots \return_k \call_1 \cdots \call_m \vAlph^*)$.
In consequence, we have:

\begin{proposition}
Consistency (resp.,  positive, negative consistency) of a ground matching-preserving rule with the language of a given \DVPA $\aut$ can be decided in polynomial time via a single reachability check on a \DVPA obtained from the product of $\aut$ with itself and a regular expression. 
\end{proposition}

\Paragraph{Matching-preserving rules with variables} 
Variables in matching-preserving rules can be eliminated as in the case of well-matched rules. 
Consider a matching-preserving rule $\lhs \to \rhs$.
It suffices to extend the alphabet with local letters $x^L, x^R$ for each variable and then consider a ground rule
obtained from  $\lhs \to \rhs$ by replacing the occurrence of every variable $X$ in $\lhs$ with $x^L$ and in $\rhs$ (if it exists)
with $x^R$.
Finally, we as in the well-matched case, we iterate over all possible assignments from variables of $\lhs \to \rhs$ to $\BalReach$. 
In consequence, we have 

\begin{theorem}
Consistency (resp., positive, negative consistency) of a matching-preserving rule with the language of a given \DVPA $\aut$ can be decided in $\coNP$.
Assuming a bound on the number of variables in rules, consistency (resp., positive, negative consistency) of a matching-preserving rule
can be check in polynomial time.
\end{theorem}


\section{Algorithms for checking consistency}
\newcommand{\BalReach}{\mathcal{B}_{(4)}}
\newcommand{\autPlus}{\aut_{\sigma}}
We develop algorithms for deciding consistency (resp., positive or negative consistency). 
We discuss separately well-matched rules (Section~\ref{s:well-matched}) and the wider class of matching-preserving rules (Section~\ref{s:matching-preserving}).
In both cases, we consider first rules without variables, and then the extension to rules with variables.
Although the asymptotic complexity is the same in both cases, the algorithms for well-matched rules are conceptually simpler and admit more efficient implementations.
Before studying particular types of rules, we make a general observation that reduces consistency checking to single-step rewriting.

\Paragraph{Reduction to single-step rewriting}
 Observe that an \SRSV $\rew$ is consistent with the language of $\aut$ if and only if for every rule $\lhs \to \rhs$ and all words $x,y$, we have
$x \lhs y \in \lang(\aut) \iff x \rhs y \in \lang(\aut)$. Indeed, if this property holds for every rule, then 
 consistency under multi-step rewriting follows by transitivity.
 This holds due to the assumption that the $\aut$ and $\rew$ are over the same alphabet~\cite{ecai25}.
%by transitivity of the equivalence 
%if $w \rewrites_{\rew} w'$, then $w \in \lang \implies w' \in \lang$.
Analogously, for positive and negative consistency it suffices to verify the property for each individual rewriting step.

\subsection{Consistency of well-matched rules}
\label{s:well-matched}

\Paragraph{Ground well-matched rules} Consider a well-matched rule $\lhs \to \rhs$ without variables, i.e., a standard string rewriting rule. 
We extend the alphabet $\vAlph$ with a fresh local letter $\#$ and construct two \DVPA over the extended alphabet.
Both \DVPA extend $\aut$ and differ only in their behavior on the letter $\#$; in $\aut_1$ for any state $q$, 
the transition over $\#$ leads to the state $q'$ such that
$\aut$ moves from $q$ to $q'$ over $\lhs$. 
Since $\lhs$ is well-matched, the run over $\lhs$ does not depend on the stack content and it does not modify the stack content. 
Thus, the effect of reading $\lhs$ in the \DVPA $\aut$ is equivalent to a transition over a single local letter. 
The \DVPA $\aut_2$ is defined analogously, but $\#$ simulates the word $\rhs$ instead. Again, word $\rhs$ is well-matched. 
Therefore, to check consistency of  $\lhs \to \rhs$  with $\lang(\aut)$ it suffices to check whether there is a word $w$ such that exactly one of \DVPA $\aut_1, \aut_2$ accepts it, 
which can be done by a single reachability check on the product \DVPA $\aut_1 \times \aut_2$, a \DVPA simulating two copies of $\aut$ in parallel.
The visibility condition ensures that a single stack over the product of stack alphabets is sufficient to simulate both \DVPA.
Checking positive (resp., negative) consistency amounts to checking whether $\lang(\aut_1) \setminus \lang(\aut_2)$  (resp., $\lang(\aut_2) \setminus \lang(\aut_1)$ for negative consistency) is non-empty, 
which again amounts to a single reachability check. 

\begin{proposition}
Consistency (resp.,  positive, negative consistency) of a ground well-matched rule with the language of a given \DVPA $\aut$ can be decided in polynomial time via a single 
reachability check on a \DVPA obtained from the product $\aut \times \aut$. 
\end{proposition}


Reachability in pushdown automata can be checked in $\bigO(|\aut|^3)$, with a slightly improved bound $\bigO(|\aut|^3/\log(|\aut|))$ for \emph{recursive state machines}~\cite{DBLP:conf/popl/Chaudhuri08}
both of which fall into $\widetilde{\bigO}(|\aut|^3)$.
Due to the lack of better upper bounds~\cite{DBLP:conf/popl/Chaudhuri08,DBLP:conf/lics/BalasubramanianCM25}, the complexity of reachability in pushdown systems is often referred to as the ``cubic bottleneck.''
Consequently, reachability in the product $\aut \times \aut$ amounts to $\widetilde{\bigO}(|\aut|^6)$. 
However, checking consistency of an SRS with the language of a \DFA can be performed locally on the structure of the \DFA without checking reachability~\cite{ecai25}. 
In contrast, the following example illustrates that for pushdown systems, checking consistency must involve reachability.

\begin{example}
Consider a \DVPA $\autB$ depicted in Figure~\ref{fig:nonsyntactic}, which is over the pushdown alphabet $(\set{c},\set{l_1, l_2}, \set{r})$ such that local letters $l_1, l_2$ induce self-loops on 
all states except for $q_0$. Its stack alphabet is a singleton $\set{a}$.
In the state $q_0$ the \DVPA moves over $l_1$ to $q_1$ and over $l_2$ to $q_2$. The automaton behaves differently in states $q_1, q_2$ if the stack is non-empty, while these states are indistinguishable 
if the stack is empty. 
Finally, $q_0$ is reachable only with non-empty stack, i.e., $q_0$ is reachable, but the configuration $(q_0, \bot)$ is not.
The rule $l_1 \to l_2$ is consistent with $\lang(\autB)$, but establishing this requires reachability analysis. 
\begin{figure}[h!]
\centering
\begin{tikzpicture}[
    shorten >=1pt, 
    node distance=3.0cm, 
    on grid, 
    auto,
    >={Stealth[bend]}
]

  % State Definitions
  \node[state,initial] (q0) {$q_I$};
  \node[state] (q1) [right=of q0] {$q_0$};
  \node[state] (q2U) [above right=of q1] {$q_1$};
  \node[state] (q2D) [below right=of q1] {$q_2$};
  \node[state] (sink) [below right=of q2U] {$q_R$};
  \node[state,accepting] (acc) [right=of q2U] {$q_F$};

  % Transition Edges
  \path[->]
    (q0) edge node [above] {$c, push(a)$} (q1)
    (q1) edge node [above] {$l_1$} (q2U)
    (q1)  edge node [above] {$l_2$} (q2D)
    (q2U) edge node [right] {$r, pop(a)$} (sink)
    (q2U) edge node [above] {$r, pop(\bot)$} (acc)
    (q2D)  edge [bend right=20] node [above left=0.6cm] {$r, pop(a)$} (sink) 
    (q2D)  edge [bend left=20] node [below right=0.6cm] {$r, pop(\bot)$} (sink);
    

    %[bend right=45]
 %   (q1)  edge [bend right=45] node [above] {$\Sigma\setminus (Y \cup \set{x})$} (qn)
    % The outgoing edge xi from q1
%    (q1)  edge [nodes={at={(1, -0.5)}}] node {$x$} ++(1.5, 0);
%    (q1)  edge node [above] {$x$} (qF)
%    (q1) edge node [left] {$Y$} (sink)
%    (q2) edge node [left] {$Y$} (sink)
%    (qn) edge node [above] {$Y$} (sink)
%    (qF) edge node [left] {$Y$} (sink)
    ;

 \draw[loop right,->] (acc)  to node[right] {$\vAlph$} (acc);
 \draw[loop right,->] (sink) to node[right] {$\vAlph$} (sink);

\end{tikzpicture}
\caption{The automaton $\autB$. All transitions other than depicted are self-loops}
\label{fig:nonsyntactic}
\end{figure}
\end{example} 


\Paragraph{Non-ground well-matched rules}
\newcommand{\autCon}{\aut_{\textrm{con}}}
\newcommand{\DeltaRule}{\Delta[\lhs \to \rhs]}
Consider a well-matched rule $\lhs \to \rhs$ with variables. 
The general approach, as in the ground case, is to reduce consistency checking to reachability in a product \VPA.
We construct a \VPA $\autCon$ that simulates two copies of $\aut$ over the alphabet extended with a fresh local letter $\#$.
The \VPA $\autCon$ is built on the product $\aut \times \aut$, whose stack track contents of both copies of $\aut$.
The fresh local letter $\#$ simulates the word $\lhs$ in the first copy and the word $\rhs$ in the second copy.
However, there are two key differences.
First, since $\lhs$ and $\rhs$ can be instantiated to (infinitely many) different words, transitions over $\#$ are non-deterministic.
Second, the transitions in each copy of $\aut$ cannot be defined independently; they must be defined jointly on the product to ensure that both sides of the rule $\lhs \to \rhs$ are instantiated to the same word. 
For example, for a rule $c X r \to X$, there may be a transition in $\aut$ from $q$ to $q_1$ over $c w_1 r$ and from $q$ to $q_2$ over $w_2$, but we must ensure that $w_1 = w_2$.
To this end, we compute a set $\DeltaRule$ of triples $(q, s_1, s_2)$ such that there are words $\alpha, \beta$ satisfying:
(1)~$s_1 \neq s_2$, 
(3)~$(q,u) \trans{\aut}{\alpha} (s_1,u)$ and $(q,u) \trans{\aut}{\beta} (s_2,u)$ for every stack content $u$, and
(4)~$\alpha \to \beta$ is an instance  of $\lhs \to \rhs$ obtained by substituting well-matched words for variables.
Using $\DeltaRule$, we define a \VPA $\autCon$ over the pushdown alphabet $\vAlph$ extended with a local letter $\#$ as an extension of $\aut \times \aut$ such that for a state $(q_1, q_2)$ of $\autCon$:
(a)~if $q_1 \neq q_2$, then $\autCon$ has no transition  over $\#$, and
(b)~if $q_1 = q_2$, then for every $(q_1, s_1, s_2) \in \DeltaRule$, 
the tuple $\langle (q_1, q_1), \#, (s_1, s_2) \rangle$ is a transition in $\autCon$. 
Intuitively, the automaton only applies the rule when both copies are in the same state, and hence words with more than one $\#$ letter are rejected.
Finally, the automaton accepts if in the reached state $(q_1, q_2)$, exactly one of $q_1, q_2$ is accepting in $\aut$. 
Observe that the \VPA $\autCon$ accepts exactly words $w$ that are witnesses of inconsistency of $\lhs \to \rhs$ with $\lang(\aut)$, i.e., 
any accepted word $w$ contains exactly one occurrence of $\#$. 
Furthermore, there is a ground instance $\alpha \to \beta$ of the rule $\lhs \to \rhs$ obtained by substituting well-matched word for variables, and 
by substituting $\#$ with $\alpha$ (resp., $\beta$) yields $w_1$ (resp., $w_2$) satisfying 
(i)~$w_1$ can be rewritten on one step to $w_2$ with the rule $\alpha \to \beta$, and 
(ii)~exactly one of $w_1, w_2$ belongs to $\lang(\aut)$.
By modifying the acceptance condition, we can also decide positive and negative consistency. 

In the following, we focus on computing $\DeltaRule$ for rules with variables. 
It suffices to compute how well-matched words transform pairs of states of the \DVPA.

\Paragraph{Computing $\DeltaRule$}
\newcommand{\lhsInst}{\bar{\alpha}}
\newcommand{\rhsInst}{\bar{\beta}}
We compute a quaternary relation $\BalReach \subseteq Q^4$ such that $(q_1, q_2, s_1, s_2) \in \BalReach$ if and only if there is a well-matched word $w$ such that
$q_1 \trans{\aut}{w} s_1$ and $q_2 \trans{\aut}{w} s_2$. 
It is computed iteratively in a similar way to reachability in pushdown automata, but applied to the product automaton $\aut \times \aut$.
 %we compute the closure with respect to 
%enclosing the word with a call and a return letter, joining a local letter and concatenation computed simultaneously.
\begin{restatable}{lemma}{BalancedReachability}
Given a \DVPA $\aut$, the relation $\BalReach$ can be computed in polynomial time in the size of $\aut$.
\end{restatable}
We extend the alphabet $\vAlph$ to $\vAlph_+$ by adding a pair of fresh local letters $x^L$ and $x^R$ for every variable $X$ occurring in $\lhs \to \rhs$.
Next, we substitute each variable $X$ with its corresponding local letters: $x^L$ for $X$ in $\lhs$ and $x^R$ for $X$ in $\rhs$, obtaining words $\lhsInst$ and $\rhsInst$ over the extended alphabet. 
Then, we iterate over all possible assignments from variables of $\lhs \to \rhs$ to $\BalReach$. 
For a given assignment $\sigma$, we extend the \DVPA $\aut$ to a \VPA $\autPlus$ working over the extended alphabet, where transitions over new letters are defined as follows.
For every variable $X$, if $\sigma(X) = (q_1, q_2, s_1, s_2)$, then
$\autPlus$ has a transition over $x^L$ from $q_1$ to $s_1$ and over $x^R$ from $q_2$ to $s_2$; there are no other transitions over $x^L$ or $x^R$.
Finally, %we evaluate $\lhs$ and $\rhs$ by evaluating $\lhsInst$ and $\beta^+$ in $\autPlus$.
for every triple $(q,s_1, s_2)$, if $(q,\bot) \trans{\autPlus}{\lhsInst} (s_1,\bot)$ and $(q,\bot) \trans{\autPlus}{\rhsInst} (s_2,\bot)$, we add the triple $(q,s_1, s_2)$  to $\DeltaRule$.
Note that the stack content $\bot$ does not influence reachability since $\lhsInst, \rhsInst$ are well-matched.


\begin{lemma}
Deciding membership in $\DeltaRule$ over well-matched rules is in \NP.
Under a fixed bound on the number of variables per rule, it is decidable in polynomial time.
\end{lemma}

In consequence, checking inconsistency can be done in \NP, which gives \coNP test for consistency:

\begin{theorem}
Consistency (resp., positive, negative consistency) of a well-matched rule with the language of a given \DVPA $\aut$ can be decided in $\coNP$.
Under a fixed bound on the number of variables per rule, it is decidable in polynomial time.
\end{theorem}


\Paragraph{Complexity of checking (positive) consistency}
The proposed algorithm checks all assignments of transitions for variables, which is exponential in the number of variables; it works in polynomial time if the number of variables is bounded. 
Significantly improving upon this complexity is unlikely as the membership problem for $\DeltaRule$ is \NP-complete in general, and 
deciding positive (resp. negative) consistency is \coNP-complete.
The hardness result holds already for \DFA, which can be considered as \DVPA over only local letters.


\begin{restatable}{theorem}{NPHardness}
(1)~The problem of deciding, given a rewrite rule $\lhs \to \rhs$, a DFA $\aut$ and its states $q, s_1, s_2$, whether $(q, s_1, s_2)$ belongs to $\DeltaRule$, is \NP-complete. 
(2)~The positive (resp. negative) consistency problem over well-matched rules against a regular language given by a \DFA is \coNP-complete.
\end{restatable}
\begin{proof}[Proof's sketch]
We show thus by reduction from the 3-SAT problem. Given a formula $\varphi$, we rename variables to obtain $\varphi^*$, in which  each variable occurs at most once. 
Then, $\aut$ moves over some instance of $\lhs$ from $q$ to $s_1$ if the variable substitution of $\lhs$ encodes a truth assignment, where all copies of the same variable have the same value.
Next, $\aut$ moves over some instance of $\rhs$ from $q$ to $s_2$  if  the variable substitution of $\rhs$ encodes a truth assignment satisfying the formula $\varphi$.
Thus, $(q, s_1, s_2)$ belongs to $\DeltaRule$ if and only if $\varphi$ is satisfiable.
Furthermore, we can reduce unsatisfiablity of $\varphi$ to positive (resp. negative) consistency. 
We ensure that whenever a truth assignment is consistent among copies of the same propositional variable, it does
not satisfy $\varphi^*$. Thus, $\varphi$ is unsatisfiable.
\end{proof}

The \coNP-hardness result for positive (resp. negative) consistency does not extend to consistency, and hence 
the complexity lower bound for checking consistency  of well-matched rules remains open.

\subsection{Consistency of matching-preserving rules}
\label{s:matching-preserving}

We turn to matching-preserving rules that are not well-matched.
The overall approach is similar to the well-matched case; we extend the alphabet with auxiliary letters and consider the product $\aut \times \aut$.
The key difference is that in well-matched rules, both sides are stack-independent and stack-neutral, whereas matching-preserving rules may depend on and modify the stack.
As a result, instead of local letters, we extend the alphabet with call and return letters. 

\Paragraph{Ground matching-preserving rules}
\newcommand{\autExt}{\aut_E}
\newcommand{\return}{\overline{\mathbf{r}}}
\newcommand{\call}{\overline{\mathbf{c}}}
Consider a matching-preserving rule $\lhs \to \rhs$ without variables that is not well-matched.
Since the rule is matching-preserving, $\lhs$ and $\rhs$ have the same number $m$ of unmatched calls and the same number $k$ of unmatched returns; 
since it is not well-matched, there are unmatched letters (at least one of $k, m$ is non-zero); assume that $k>0$ (the case $k=0$ and $m>0$ is similar). 
Furthermore,  $\lhs$ and $\rhs$ can be partitioned into well-matched words and unmatched calls and returns, where
unmatched returns precede unmatched calls.
It follows that  
\begin{align*}
&&\lhs &= v_0^L \cdot r_1^L \cdot v_1^L \ldots r_k^L \cdot v_k^L \cdot c_1^L \cdot v_{k+1}^L \cdot c_2^L \ldots c_m^L \cdot v_{k+m}^L \\
&&\rhs &= v_0^R \cdot r_1^R \cdot v_1^R \ldots r_k^R \cdot v_k^R \cdot c_1^R \cdot v_{k+1}^R \cdot c_2^R \ldots c_m^R \cdot v_{k+m}^R
\end{align*}
where $v_i^L, v_i^R$ are well-matched words, $r_i^L, r_i^R$ are return letters and $c_i^L,c_i^R$ are call letters. 
As in the well-matched case, let $\autExt$ be an extension of the product automaton $\aut \times \aut$, i.e., 
its set of states is $Q \times Q$ and its stack alphabet is $\stackAlphabet \times \stackAlphabet$.
The automaton $\autExt$ works over an extended alphabet, where 
the additional letters are $k$ fresh return letters
$\return_1, \ldots, \return_k$
and $m$ fresh call letters $\call_1, \ldots, \call_m$. 
The transition over $\return_1$ simultaneously  simulates reading $v_0^L r_1^L v_1^L$  on the first component of $\autExt$ and $v_0^R r_1^R v_1^R$ on the second.
In doing so, it pops one symbol from the product stack, which encodes a pair of stack symbols from the two copies of $\aut$.
Formally, $\delta_E((q_1, q_2),\return_i, (b_1, b_2)) = (s_1, s_2)$ if $(q_1,u b_1) \trans{\aut}{v_0^L r_1^L v_1^L} (s_1, u)$ and 
$(q_2,u b_2) \trans{\aut}{v_0^R r_1^R v_1^R} (s_2, u)$ for every stack content $u$. 
For the remaining return letters $\return_i$, where $i \in \set{2, \ldots,k}$,
the transitions on $\return_i$ simulate reading  $r_i^L v_i^L$  on the first component and $r_i^R v_i^R$ on the second component.
For each new call letter $\call_i$, the transition over $\call_i$ simultaneously  simulates reading $c_i^L v_{k+i}^L$ on the first component and $c_i^R v_{i+k}^R$ on the second.
Finally, a state $(q_1, q_2)$ of $\autExt$ is accepting if exactly one of $q_1, q_2$ is accepting in $\aut$.

The following condition holds:
\begin{claim}
A  word  $w_1 \return_1 \ldots \return_k \call_1 \cdots \call_m w_2$, where $w_1, w_2$ are over the original alphabet $\vAlph$, is accepted by $\autExt$  
if and only if $\aut$ accepts exactly one of the words $w_1 \lhs w_2, w_1 \rhs w_2$, which holds precisely when $\lhs \to \rhs$ is inconsistent with $\lang(\aut)$.
\end{claim}
Therefore, the rule $\lhs \to \rhs$ is consistent with $\lang(\aut)$ if and only if
the language $\lang(\autExt) \cap \lang(\vAlph^* \return_1 \ldots \return_k \call_1 \cdots \call_m \vAlph^*)$ is empty.

By changing the accepting states of $\autExt$ as discussed in the well-matched case, we can decide positive and negative consistency 
using emptiness of $\lang(\autExt) \cap \lang(\vAlph^* \return_1 \ldots \return_k \call_1 \cdots \call_m \vAlph^*)$.
In consequence, we have:

\begin{proposition}
Consistency (resp.,  positive, negative consistency) of a ground matching-preserving rule with the language of a given \DVPA $\aut$ can be decided in polynomial time via a single reachability check on a \DVPA obtained from the product $\aut \times \aut \times \autB$, where $\autB$ is a \DFA depending only on the rule.
\end{proposition}

Now, we can move to rules with variables. We only briefly comment on the extension as it is similar as for well-matched rules.

\Paragraph{Matching-preserving rules with variables} 
Variables in matching-preserving rules can be eliminated as in the case of well-matched rules. 
Consider a matching-preserving rule $\lhs \to \rhs$.
It suffices to extend the alphabet with local letters $x^L, x^R$ for each variable and then consider a ground rule
obtained from  $\lhs \to \rhs$ by replacing the occurrence of every variable $X$ with $x^L$ in $\lhs$ and with $x^R$ in $\rhs$ whenever appears.
Finally, as described in Section~\ref{s:well-matched} for the well-matched case, we iterate over all possible assignments from variables of $\lhs \to \rhs$ to $\BalReach$. 
In consequence, we have the following:

\begin{theorem}
Consistency (resp., positive, negative consistency) of a matching-preserving rule with the language of a given \DVPA $\aut$ can be decided in $\coNP$.
Under a fixed bound on the number of variables per rule, it is decidable in polynomial time.
\end{theorem}




%%% IMPROVMENTS

%The complexity discussion is split across two locations in the well-matched subsection. The proposition on ground rules states polynomial-time decidability, then the cubic bottleneck discussion follows, then the example, then the non-ground case introduces NP membership, and finally the hardness results appear at the end. This means the reader encounters upper bounds, then a digression, then more upper bounds, then lower bounds — all interleaved with algorithmic content. Consider consolidating the complexity discussion: state all upper bounds first (ground: polynomial, non-ground: coNP / polynomial with bounded variables), then present the hardness results together at the end of the subsection. This would give the reader a cleaner picture.

%The example is well-motivated but interrupts the algorithmic narrative. It appears between the ground and non-ground cases, which is a natural break point. However, it makes the point that reachability is necessary, which is a lower-bound-style observation. Placing it after the complexity results — rather than between the two algorithmic cases — would keep the algorithmic development uninterrupted and group all complexity-related discussion together.

%The matching-preserving subsection feels noticeably shorter and less developed than the well-matched one. The ground case is explained in reasonable detail, but the treatment of variables is compressed into two sentences that simply say "do the same as before." While this is accurate, it leaves the reader unsure whether there are genuinely no new complications. If the argument is indeed identical, a single explicit sentence confirming this — "The construction for variables in matching-preserving rules introduces no new difficulties beyond those already handled in the well-matched case" — would be more reassuring than a terse forward reference. If there are subtle differences, they should be flagged.

%The two subsections are not fully parallel in structure. The well-matched subsection has: ground rules → proposition → complexity discussion → example → non-ground rules → \DeltaRule computation → lemma → theorem → hardness discussion. The matching-preserving subsection has: ground rules → proposition → variables in two sentences → theorem. The asymmetry is partly justified by the fact that the variable elimination is identical, but the reader may feel that something is missing. Either making the structure more explicitly parallel, or acknowledging the asymmetry and explaining why it is justified, would improve the overall coherence.

% The open problem at the end of the well-matched subsection is a reasonable closing remark, but it lands somewhat flatly. Since it concerns only the lower bound for full consistency (not positive/negative), it is a precise and meaningful observation. Consider framing it slightly more prominently — for instance, as a named open question — so it does not read as an afterthought.

\section{Rule synthesis}
In the previous section, 
we have studied how advice provided via an \SRS reduces the number of equivalence queries necessary to learn an unknown language.
In this section, we address the converse: \emph{synthesis} of rewrite rules consistent with a given visibly pushdown language. 
Synthesized rewrite rules can serve as an explanation of the visibly pushdown language under consideration, capturing its key properties through concise rewrite rules.

There are two possible variants of this problem:
\begin{description}
\item[\textit{Synthesis from Automata}:] Given a \DVPA, generate string rewrite rules consistent with the language of the given \DVPA.
\item[\textit{Synthesis from Oracles}:] Given an oracle for an unknown visibly pushdown language, generate string rewrite rules consistent with the language. 
\end{description}

While \textit{Synthesis from Oracles} seems more natural, we show that even given a \DVPA, the synthesis problem is challenging for two reasons: 
high computational complexity and the lack of a clear objective --- it is not obvious what constitutes a ``relevant'' rule. 
Therefore, \textit{Synthesis from Oracles} suffers from these same shortcomings, and further progress in \textit{Synthesis from Automata} is required before \textit{Synthesis from Oracles} can be effectively addressed.
Consequently, for the remainder of this section we focus on \textit{Synthesis from Automata}.



We establish a polynomial upper bound on the lengths of words in minimal-length rules consistent with $\aut$, which implies that such rules
can be computed in non-deterministic polynomial time. 
In order to achieve that, we give  sufficient conditions for a pair of words to form a rewrite rule consistent with $\lang(\aut)$.
We consider first rules without variables, and then move to rules with variables.

We consider a \DVPA $\aut$ with the set of states $Q$ over a non-trivial alphabet, i.e., either there are at least two local letters or at least one pair of a call and a return letter.
Under these assumptions, there are exponentially many well-matched words.

\Paragraph{The length of a minimal consistent well-matched rule: the ground case}
Consider a well-matched word $w$. We associate with this word a function $f_w \colon Q \to Q$.
Such functions correspond to equivalence classes of the syntactic congruence of a \DFA.
Formally, for a state $q \in Q$ we define $f_w(q) = q'$ if and only if for some stack content $u$, we have $(q,u) \trans{\aut}{a} (q',u)$.
Note that the existential quantifier over $u$ can be replaced by a universal one, since transitions over well-matched words are stack-independent. 
If $w_1, w_2$ are well-matched words such that $f_{w_1} = f_{w_2}$, 
then for all words $x,y$ we have $x w_1 y \in \lang(\aut) \iff x w_2 y \in \lang(\aut)$~\cite[Theorem~1]{alur2005congruences}.
It follows that $w_1 \to w_2$ is consistent with $\lang(\aut)$. 
Setting $n = |Q|$, there are $n^n$ functions from $Q$ to $Q$, and hence there are $n^n = 2^{\bigO(n \log n)}$ distinct functions $f_w$ over all well-matched words $w$.
There are $2^{\Theta(k)}$ well-matched words of the length $k$, and hence among well-matched words of length $\bigO(n \log n)$ 
there is a pair of well-matched words $\lhs, \rhs$ such that $f_{\lhs} = f_{\rhs}$. 
It follows that $\lhs \to \rhs$ is a ground well-matched rule consistent with $\lang(\aut)$. 

\Paragraph{The length of a minimal consistent well-matched rule: the general case}
For rules with variables, we consider only a subset of rules and only a sufficient condition for consistency. 
Nevertheless, we obtain asymptotically the same bound. 
Consider a rewrite rule with variables $\lhs \to \rhs$ of the form:
\begin{align*}
 && w_0 X_1 w_1 X_2 \ldots  X_k w_{k} &\to v_0 X_1 v_1 X_2 \ldots  X_k v_{k}
\end{align*}
where the order of variables in $\lhs$ and $\rhs$ is the same and all words $w_i,v_i$ are well-matched words over $\vAlph$, i.e.,
do not contain variables.
Observe that if the corresponding words $w_i$ and $v_i$ define the same function, i.e., 
$f_{w_0} = f_{v_0}, \ldots, f_{w_{k}} = f_{v_k}$, then the rewrite rule $\lhs \to \rhs$ is consistent with the language of $\aut$. 
To see this, consider an instantiation of $\lhs \to \rhs$, i.e., 
consider a sequence of well-matched words $b_1, \ldots, b_k$, which we substitute for $X_1, \ldots, X_k$
obtaining the ground rewrite rule:
\begin{align*}
&&	w_0 b_1 w_1 b_2 \ldots  b_k w_{k} &\to v_0 b_1 v_1 b_2 \ldots  b_k v_{k}
\end{align*}
Observe that the left hand side and the right hand side define the same function, and hence this ground rule is 
consistent with $\lang(\aut)$.
It follows that every rule from $\sem{\lhs \to \rhs}$ is consistent with $\lang(\aut)$, which by definition means that 
$\lhs \to \rhs$ is consistent with $\lang(\aut)$.
Now, setting $n = |Q|$, there are $n^{k\cdot n}$ different combinations of $k$ functions from $Q$ to $Q$. 
The number of well-matched words is $2^{\Omega(n)}$ and hence there is a constant $c>0$ such that
among words of the form $w_0 X_1 w_1 X_2 \ldots  X_k w_{k}$ such that 
the sum of lengths of $w_0, \ldots, w_k$ is bounded by  ${c \cdot k\cdot n \cdot \log(n)}$, 
there is a pair of words that forms a rewriting rule consistent with $\lang(\aut)$.

%\Paragraph{Rule synthesis in non-deterministic polynomial time}
The counting argument shows that for a given $\DVPA$ $\aut$ there are short (of polynomial length in $|\aut|$) 
rewrite rules that are consistent with $\lang(\aut)$. 
Since testing whether a given rule (with a bounded number of variables)
is consistent with $\lang(\aut)$ is decidable in polynomial time, we can find (some) rules consistent with $\lang(\aut)$ in non-deterministic polynomial time. 

\Paragraph{Rules of exponential length}
While minimal-length rules are of polynomial length, some consistent rules are of exponential length even if we consider ``nonredundant'' rules. 
Even in the \DFA case, consider automata $\aut_n$ such that $|Q| = n$  and for all $a \in \Sigma$, the function $f_a$ is the same permutation of states of an exponential order, e.g., the composition 
of disjoint cycles of prime orders $p_1, \ldots, p_k$ has the order $p_1 \cdot \ldots \cdot p_k$.
Then, there exist words $w \neq \epsilon$ such that $w \to \epsilon$ is a rule consistent with $\lang(\aut_n)$, but every such word $w$ has length at least exponential in $n$.

While we have shown that rules can be generated in non-deterministic polynomial time, 
a natural question is whether this is optimal.

\Paragraph{Complexity of finding consistent rules}
We have established an upper bound on the minimal length of words defining the same function. 
We show that computing the exact minimal length of such words is difficult, which implies that
finding a minimal-length rewriting rule, which is consistent with a given \DVPA is difficult as well.
Furthermore, the problem is already difficult  for \DFA.

We consider the \emph{minimal collision} problem defined as follows: given a \DFA $\aut$ and $k>0$, 
does there exist a pair of different words $w_1, w_2$ such that (1)~$|w_1|,|w_2| \leq k$, and
(2)~$f_{w_1} = f_{w_2}$. As we have discussed above for the ground rewrite rules, if $|\Sigma| \geq 2$, then for every $k > |\aut|\cdot \log(|\aut|)$ the answer is yes. 
Consequently, the problem is in \NP. We show that it is \NP-complete.

\begin{restatable}{lemma}{NPHardnessOfClash}
The minimal collision problem is \NP-complete.
\end{restatable}
\begin{proof}[Proof's sketch]
The proof is by reduction from the 3-SAT problem.
The main difficulty is to ensure that if the formula is unsatisfiable, 
then no pair of short words defines the same function from $Q$ to $Q$ .
We construct a gadget \DFA such that unless one of the designed letters occurs in a word, 
each word of length at most $n$ generates a unique function.
We consider the alphabet consisting of literals from a given 3-CNF formula.
With that gadget \DFA, we enforce that short words having a collision must contain a variable or its negation, i.e., 
represent an assignment, and these words must contain a literal from each clause.
\end{proof}


% In the minimal collision problem, both sides of the rule are unknown words.
% We discuss another decision problem associated with rule synthesis, in which the right-hand side of the rule is fixed. 
% For simplicity, we consider the right-hand side to be a local letter.
% The \emph{bounded rewriting to single letter} problem asks, given a \DFA $\aut$ over $\Sigma$,  a letter $a \in \Sigma$, and a bound $k>0$, 
% whether there exists a word $w \in (\Sigma \setminus \set{a})^*$ such that 
% (1)~$w \to a$ is a rule consistent with $\lang(\aut)$, and
% (2)~$|w| \leq k$. 
% This problem is closely related to the \emph{membership}  problem for monoids~\cite{Kozen77} and groups~\cite{DBLP:journals/tcs/Jerrum85}.
% The membership problem for monoids reduces to the \emph{bounded automata intersection} problem~\cite{Kozen77,DBLP:conf/wia/Wareham00}, which 
% asks, given $m,k>0$ and \DFA $\aut_1, \ldots, \aut_m$, whether there exists a word $w$ such that $|w| \leq k$ and $w \in \lang(\aut_i)$ for every $i$.
% The bounded automata intersection problem is \PSPACE-complete if $k$ is given in binary~\cite{Kozen77} and \NP-complete if $k$ is given in unary~\cite{DBLP:conf/wia/Wareham00}.
% The same complexity bounds apply to the bounded rewriting to single letter problem as there is a straightforward reduction from the bounded automata intersection problem.
% Clearly, the problem is in $\PSPACE$ for $k$ given in binary and $\NP$ for $k$ given in unary.

% \begin{restatable}{lemma}{ReductionToLocal}
% The bounded rewriting to a single letter with the parameter $k$ is 
% (a)~\PSPACE-complete, if $k$ is given in binary, and
% (b)~\NP-complete, if $k$ is given in unary.
% \end{restatable}

In summary, there is no algorithm for generating minimal-length consistent rules in deterministic polynomial time unless \PTIME = \NP.
%\todo{ soften it slightly}
%Consequently, further research on the rule synthesis problem should focus on efficient heuristics rather than exact algorithms.

%% \cite{DBLP:journals/aaecc/ThiemannZGS08}
%% Differences between TRS and SRS and why synchornzing words do not reduce to SRS rule synthesis
%% with TRS we can express that a rule correwsponds to a reset word. It is not possible with an SRS.
%% 


\section{Empirical evaluation}
We evaluate our framework by expressing natural properties of \VPLs as rewrite rules with variables. 
Our experiments show that learning with advice yields 
a substantial reduction in the number of equivalence queries, runtime reduction and improved accuracy,
demonstrating the practical advantages of the proposed method.    

\subsection{Examples of rewrite rules}
We consider the following rules, which express natural properties of \VPLs.
\begin{itemize}
\item \emph{(Well-matched) commutativity}: $c_1 c_2 X r_2 r_1 \to c_2 c_1 X r_1 r_2$. Consider a \DVPA 
modelling a system consisting of two pushdown components such that $c_1, r_1$ are actions of the first component,
while $c_2, r_2$ are actions of the second. Then, the trace language of the whole system satisfies 
(well-matched) commutativity.
\item \emph{(Well-matched) cancellation}: $c X r \to \epsilon$. Consider a markup language, in which $c$ is a tag opening comment section,
while $r$ is a tag closing it. In such a case, (well-matched) cancellation holds.
\item \emph{(Well-matched) idempotency}: $c c X r r \to c X r$. Consider a markup language, in which $c$ is a tag activating certain mode 
(e.g. bold text) and $r$ is a tag deactivating it. In this setting (well-matched) idempotency holds.
\end{itemize}

% reserch questions

\subsection{Implementation}

\Paragraph{Baseline}
As a baseline for our experiments, we have implemented a polynomial-time active learning algorithm for \DVPA~\cite{mfcs22}.
The algorithm from~\cite{mfcs22} assumes a stronger teacher model than the typical MAT; it allows extended membership queries
that both read and modify the stack content. Moreover, these queries depend on a concrete \DVPA rather than \VPL, and hence
in this framework the automaton is learned rather than the language. This allowed for
 the learning algorithm to work in polynomial time. 

\Paragraph{Equivalence oracle}
We have implemented two types of equivalence test: the exact test based on reachability in the product automaton, and
the probabilistic test based on testing conformance of automata on a large number ($10^5$ and $10^6$) of randomly generated words
of a random length bounded by $100$.
The probabilistic equivalence test is significantly faster than the exact one (from 10x to even 100x faster for 25-state automata).

\Paragraph{The advice mechanism}
We have extended the baseline algorithm with the advice mechanism, which intercepts equivalence queries before
passing them to the equivalence oracle, and performs the consistency check.
The advice is given via rewrite rules with a {single} variable.
We have implemented the exact algorithm presented in the paper, and two independent heuristics. 
The first heuristic is to instantiate a rule with variables to 40--200 ground rules and work with those.
%We select well-matched words to be pairwise maximally different functions from $Q$ to $Q$, using them as candidates for substitution.
The second heuristic is the probabilistic consistency test, which, analogously to the equivalence test, is
approximated by evaluation of a large number of randomly generated words.
 

\subsection{Experiments}

\Paragraph{Benchmarks}
We have generated random benchmark \DVPA with multiple parameters such as the size of the alphabet, 
the size of the stack alphabet, number of states and others.
consistent with each of the advice rules discussed above. For instance, to randomly generate 
a \DVPA consistent with {(well-matched) commutativity}, we generate two random \DVPA over disjoint alphabets, take the product of these automata, where undefined transitions are defined as self-loops. 

We have conducted a series of experiments to address the following research questions:
\begin{enumerate}[Q1]
\item What is the average reduction in the number of equivalence queries avoided using advice? 
\item Does the use of the advice mechanism lead to better running time of the whole learning setup 
(i.e., both the learning algorithm and the equivalence oracle)? 
\item Does the use of the advice mechanism lead to higher reliability in presence of
the probabilistic equivalence oracle?
\end{enumerate}

Research questions Q1, Q2 have been investigated in the \textbf{exact setting}, where the equivalence check as well as consistency check are 
computed based on  reachability, while questions Q2 and Q3 have been investigated in the \emph{probabilistic setting}, where 
the equivalence and consistency checks are approximated.



\Paragraph{The exact setting}
The exact equivalence checks dominate the overall computation time. 
When no inconsistency with an \SRS is detected, we perform the actual equivalence check, which involves two reachability checks.
However, the consistency check is computationally cheaper  because the product of the candidate automaton with itself is considerably smaller than the product of the candidate automaton with the target automaton.
The results are aggregated over $100$ randomly generated automata for each state count.
The number of well-matched words used as substitutions for the rule variable was limited to $40$.
We observe that for the cancellation property ($c X r \to \epsilon$), increasing
the number of substitutions does not improve results; the substitution $X \mapsto \epsilon$ was sufficient. 
However, this is not true for reliability, which is discussed below.


\begin{figure}[t]
\centering

\begin{subfigure}{0.48\linewidth}
\centering
\begin{tikzpicture}
\begin{axis}[
    width=\linewidth,
    height=5cm,
    xlabel={Number of states},
    ylabel={Reduction in queries [\%]},
    legend to name=sharedlegend,
    legend columns=3,
    grid=both
]

% === Cancellation ===
\addplot coordinates {
    (4, 37.5)
    (6, 40.9)
    (8, 41.9)
    (10, 45.7)
    (12, 47.1)
    (14, 54.3)
    (16, 57.0)
    (18, 55.9)
    (20, 51.7)
    (22, 52.7)
    (24, 52.2)
    (26, 55.5)
    (28, 64.0)
    (30, 62.2)
};
\addlegendentry{Cancellation}

% === Commutativity ===
\addplot coordinates {
    (6, 26.3)
    (9, 14.6)
    (12, 30.9)
    (15, 25.0)
    (16, 26.2)
    (18, 13.0)
    (20, 32.7)
    (30, 22.3)
};
\addlegendentry{Commutativity}

% === Idempotency ===
\addplot coordinates {
    (4, 6.7)
    (6, 17.4)
    (8, 9.1)
    (10, 25.0)
    (12, 18.6)
    (14, 15.0)
    (16, 21.4)
    (18, 23.5)
    (20, 18.3)
    (22, 11.7)
    (24, 25.6)
    (26, 17.4)
};
\addlegendentry{Idempotency}

\end{axis}
\end{tikzpicture}
\caption{Reduction in the number of queries}
\end{subfigure}
\hfill
\begin{subfigure}{0.48\linewidth}
\centering
\begin{tikzpicture}
\begin{axis}[
    width=\linewidth,
    height=5cm,
    xlabel={Number of states},
    ylabel={Execution time change [\%]},
    grid=both
]

% === Cancellation ===
\addplot coordinates {
    (4, -58.4)
    (6, -55.5)
    (8, -52.3)
    (10, -52.7)
    (12, -45.3)
    (14, -31.5)
    (16, -3.0)
    (18, -6.4)
    (20, -10.8)
    (22, 6.5)
    (24, 8.7)
    (26, 2.0)
    (28, 21.8)
    (30, 8.6)
};

% === Commutativity ===
\addplot coordinates {
    (6, -84.9)
    (9, -78.1)
    (12, -60.2)
    (15, -51.5)
    (16, -17.1)
    (18, -72.1)
    (20, -63.0)
    (30, 11.1)
};

% === Idempotency ===
\addplot coordinates {
    (4, -89.7)
    (6, -76.2)
    (8, -83.0)
    (10, -49.4)
    (12, -46.1)
    (14, -32.7)
    (16, -31.6)
    (18, -18.6)
    (20, -33.2)
    (22, -44.1)
    (24, 5.1)
    (26, -44.6)
};

\end{axis}
\end{tikzpicture}
\caption{Relative change in execution time}
\end{subfigure}

\vspace{2mm}
\ref{sharedlegend}

\caption{Impact of SRS on the number of equivalence queries and execution time across different SRS rule types}

\end{figure}

\FloatBarrier

\Paragraph{The probabilistic setting}
%In this section, we evaluate two aspects: the impact of SRS on reliability and the reduction in processing time.
Since the \SRS enables equivalence testing to be performed in two stages, 
we compare the algorithm with advice against the baseline, using two configurations for the latter: 
$10^5$ and $2 \cdot 10^5$ random words.
The second configuration checks whether the improvements observed with advice are genuinely due to the \SRS, 
rather than a result of evaluating a larger number of random words.
The experimental setup consists of $10^4$ randomly generated automata with $10$--$25$ states.
%The input alphabet was generated individually for each automaton, with slight variations depending on the considered property.
%In particular, the number of symbols of each type ranged from $2$ to $5$, while the number of stack symbols ranged from $2$ to $16$.
The number of well-matched words used for substitutions was limited to $200$.

\begin{table}[H]
\centering
\begin{tabular}{cc|ccc}
\hline
SRS rule & Metric
& $10^5$ random words
& $2\cdot10^5$ random words
& $10^5$ with SRS \\
\hline
\multirow{2}{*}{Cancellation}
& Success ratio & 30.29\% & 40.1\% & 57.19\% \\
& Avg. time [s] & 0.69 & 1.28 & 2.6 \\
\hline
\multirow{2}{*}{Commutativity}
& Success ratio & 73.93\% & 78.59\% & 74.11\% \\
& Avg. time [s] & 0.69 & 0.89 & 2.0 \\
\hline
\multirow{2}{*}{Idempotency}
& Success ratio & 91.0\% & 93.24\% & 91.69\% \\
& Avg. time [s] & 0.38 & 0.78 & 0.92 \\
\hline
\end{tabular}
\caption{Comparison of success ratio and average execution time for different SRS rule types}
\end{table}

\Paragraph{Summary}
Our results show a reduction of approximately $20\%$ in the equivalence queries for commutativity and idempotency,
and around $50\%$ for cancellation rules. 
While time measurements indicate that the \SRS generally increases execution time, it reduces it for the largest instances, suggesting the benefit grows with automaton size.
Based on prior experiments for \DFA~\cite{ecai25}, we expected significant gains for commutativity; however, the results differ.
We suspect the generated \DVPA consistent with commutativity were too simple for the advice to significantly impact performance.
Reliability results did not show a dramatic improvement, but the \SRS did improve the success ratio.
The overhead mainly stems from the auxiliary heuristic for selecting well-matched words. 
While this cost is negligible compared to an exact check, it is noticeable in the probabilistic setting. 
However, if word evaluation becomes more expensive (e.g., using an LLM to evaluate words instead of \DVPA), this overhead vanishes, suggesting \SRS is a valuable feature for heuristic approaches.
 


%\section{Experiments}
%\input{experiments}



\bibliographystyle{plain}
\bibliography{papers}

\NPHardness*
\begin{proof}

\end{proof}


\NPHardnessOfClash*
\begin{proof}
Let us fix a parameter $n$. Before the reduction, we construct a gadget that distinguishes all words of length at most $n$ unless they consist one of the designated letters. 
Let $Y \subseteq \Sigma$ and $x \in \Sigma \setminus Y$. Consider the automaton $\autB_{Y,x}$ depicted below:

\begin{figure}[h!]
\centering
\begin{tikzpicture}[
    shorten >=1pt, 
    node distance=2.5cm, 
    on grid, 
    auto,
    >={Stealth[bend]}
]

  % State Definitions
  \node[state] (qn) {$q_n$};
   \node[draw=none] (dots) [right=of qn] {$\dots$};
  \node[state] (q2) [right=of dots] {$q_2$};
  \node[state] (q1) [right=of q2] {$q_1$};
\node[state] (qF) [right=of q1] {$q_F$};

\node[state] (sink) [below=of q2] {$s$};

  % Transition Edges
  \path[->]
    (qn) edge node [above] {$\Sigma\setminus Y$} (dots)
    (dots) edge node [above] {$\Sigma\setminus Y$} (q2)
    (q2)  edge node [above] {$\Sigma\setminus Y$} (q1)
    % The top arc from qn back to q1
    (q1)  edge [bend right=45] node [above] {$\Sigma\setminus (Y \cup \set{x})$} (qn)
    % The outgoing edge xi from q1
%    (q1)  edge [nodes={at={(1, -0.5)}}] node {$x$} ++(1.5, 0);
    (q1)  edge node [above] {$x$} (qF)
    (q1) edge node [left] {$Y$} (sink)
    (q2) edge node [left] {$Y$} (sink)
    (qn) edge node [above] {$Y$} (sink)
    (qF) edge node [left] {$Y$} (sink)
    ;

 \draw[loop,->] (qF)  to node[above] {$\Sigma \setminus Y$} (qF);
 \draw[loop right,->] (sink) to node[right] {$\Sigma$} (sink);

\end{tikzpicture}
%\label{The gadget \DFA $\autB_{Y,x}$, which parially distingushes words of length at most $n$, which do not contain any letter from $Y$ }
\end{figure}

Observe that for all words $w_1, w_2$ of length bounded by $n$, if $f^{\autB_{Y,x}}_{w_1} = f^{\autB_{Y,x}}_{w_2}$, then one of the following holds:
\begin{itemize}
    \item each of the words $w_1, w_2$ contains at least one letter from $Y$; in that case $f_{w_1}, f_{w_2}$ are constant functions equal 
    $s$ for every argument,
    \item the sets positions at which $x$ occurs in $w_1$ and resp. $w_2$ are equal; 
    indeed, for any $w \in (\Sigma \setminus Y)^{\leq n}$ we have
    $f_{w}(q_i) = q_F$ if and only if $w[i] = x$.
\end{itemize}
Now, consider $\autB_Y$ being the disjoint union of automata $\autB_{Y,x}$ over all $x \in \Sigma\setminus Y$.

Observe that for all words $w_1, w_2$ of length bounded by $n$, if $f^{\autB_{Y}}_{w_1} = f^{\autB_{Y}}_{w_2}$, then one of the following holds:
\begin{itemize}
    \item each of the words $w_1, w_2$ contains at least one letter from $Y$; in that case $f_{w_1}, f_{w_2}$ are constant functions equal 
    $s$ for every argument,
    \item words $w_1$ and $w_2$ are equal.
\end{itemize}

Having the automaton $\autB_Y$ we reduce the 3-SAT problem to the minimal collision problem.
Let $\varphi$ be a 3-CNF formula over variables $x_1, \ldots, x_n$.
Consider $\Sigma = \{x_1, \bar{x_1}, \ldots, x_n, \bar{x_n}, a, b \}$.
With each clause $c_j$ we associate the subset $C_j$ of $\Sigma$ consisting of the literals from $c_j$.
Furthermore, for every variable $x_i$ we define the set $X_i = \set{x_i, \bar{x_i}}$.
Now, consider a \DFA $\aut_0$ being the disjoint union of automata $\autB_{C_1}, \ldots, \autB_{C_m}$ over all clauses of $\varphi$ and
automata $\autB_{X_1}, \ldots, \autB_{X_n}$ over all variables of $\varphi$. 
We extend $\aut_0$ to $\aut$ with additional states such that over these states letters from $\Sigma \setminus \set{a,b}$ are self loops, and
letters $a,b$ ensure that every component of the disjoint union is reachable. 
Let $k$ in the instance of the minimal collision problem be equal $n$, the number of variables.

Now, consider a pair of different words $w_1, w_2$ such that $f^{\aut}_{w_1} = f^{\aut}_{w_2}$ and $|w_1|,|w_2| \leq n$.
Then, for every automaton $\autB_{X_1}, \ldots, \autB_{X_n}$, since $w_1 \neq w_2$, each word contains at least one letter from each set $X_i$.
There are $n$ such sets, which are disjoint, and hence $w_1$ and $w_2$ contain exactly one letter from each of the sets $X_i$, i.e., 
$w_1$ and $w_2$ represent consistent variable assignments (they either contain $x_i$ or $\bar{x_i}$).
Furthermore, for every clause $c_j$ each word $w_1, w_2$ contains at least one letter corresponding to some literal from $c_j$.
Therefore, $w_1$ and $w_2$ represent (possibly different) satisfying assignments for $\varphi$.

Conversely, assume that $\varphi$ has a satisfying assignment. 
Consider two words $w_1, w_2$ representing this variable assignment with a different order or literals.
Then, for every component $\autB_{Y,x}$ of $\aut_0$, the functions $f_{w_1}^{\autB_{Y,x}}, f_{w_2}^{\autB_{Y,x}}$ are constant functions equal $s$.
Furthermore, on states in $\aut$, which are not in $\aut_0$, both functions $f_{w_1}^{\aut}, f_{w_2}^{\aut}$ coincide with the identity.
Consequently, we have
(a)~$w_1 \neq w_2$, (b)~$|w_1| = |w_2| = n$, and
(c)~$f_{w_1}^{\aut_0}, f_{w_2}^{\aut_0}$, i.e., 
the constructed instance of the minimal collision problem is positive.
\end{proof}

\end{document}
